\subsection{模的同型分量}

定义 4. —— 设 A 是一个环，M 是一个 A-模，S 是一个单 A-模。M 的 S 型同型分量记作 \(M_{S}\)，它是 M 中同构于 S 的子模之和。

显然，\(M_{S}\) 是 M 的 S 型同型的最大子模。由于 \(M_{S}\) 的每个子模都是 S 型同型模（VIII, 第 61 页，命题 2），故对 M 的每个子模 N 都有 \(N_{S} = M_{S} \cap N\)。

若 \(S'\) 是同构于 S 的单 A-模，则显然有 \(M_{S} = M_{S'}\)，因而 \(M_{S}\) 只依赖于 S 的类（VIII, 第 51 页）。

设 M 是一个 A-模。M 存在一个最大的半单子模，称为 M 的底座；它是 M 的诸单子模之和，也是 M 的诸同型分量之和。特别地，M 是半单模当且仅当它等于自身的底座。

命题 4. —— 设 A 是一个环。用 S 表示单 A-模的类组成的集合。设 M 是一个半单 A-模。

\begin{enumerate}
\def\labelenumi{\alph{enumi})}
\item
  模 M 是其同型分量族 \((\mathrm{M}_{\lambda})_{\lambda\in\mathcal{S}}\) 的直和。
\item
  设 M 是单子模族 \((\mathrm{N}_{i})_{i\in\mathrm{I}}\) 的直和。对每个 \(\lambda\in\mathscr{S}\)，用 \(\mathrm{I}(\lambda)\) 表示这样的指标 \(i\in I\) 组成的集合：\(N_{i}\) 属于类 \(\lambda\)。我们有 \(M_{\lambda}=\bigoplus_{i\in\mathrm{I}(\lambda)}N_{i}\)。
\item
  若 M 是有限生成的，则诸 \(\lambda\in S\) 中使 \(M_{\lambda}\neq0\) 者组成的集合是有限的。
\item
  设 \(\mathrm{N}\) 是 \(\mathrm{M}\) 的任意一个子模。对每个 \(\lambda \in \mathcal{S}\)，我们有 \(\mathrm{N}_{\lambda} = \mathrm{N} \cap \mathrm{M}_{\lambda}\) 与 \((\mathrm{M} / \mathrm{N})_{\lambda} = (\mathrm{M}_{\lambda} + \mathrm{N}) / \mathrm{N}\)。
\end{enumerate}

由于 M 是半单的，它是族 \((\mathrm{M}_{\lambda})_{\lambda\in\mathcal{S}}\) 之和；我们来证明这个和是直和。设 \(\lambda\in S\)。把族 \((\mathrm{M}_{\mu})_{\mu\in\mathcal{S}-\{\lambda\}}\) 之和记为 \(M_{\lambda}^{\prime}\)。模 \(M_{\lambda}^{\prime}\) 是一族不同构于 \(\lambda\) 的单模的直和（VIII, 第 56 页，定理 1）。由 VIII, 第 56 页定理 1 的推论 3，\(M_{\lambda}^{\prime}\) 不含任何类为 \(\lambda\) 的单子模。因此我们有 \(M_{\lambda}\cap M_{\lambda}^{\prime}=0\)。断言 a) 于是得证。按构造，我们有 \(M_{\lambda}\supset\bigoplus_{i\in\mathrm{I}(\lambda)}N_{i}\)，故断言 b) 由 II, §1, No.~8, 第 208 页的注 1 得出。

断言 c) 由 a) 及 II, §1, No.~12, 第 221 页的命题 23 得出。

设 N 是 M 的一个子模，\(\lambda\in S\)。N 的同型分量 \(N_{\lambda}\)（\(\lambda\) 型）包含于 \(M_{\lambda}\)，且 \(M_{\lambda}\cap N\subset N_{\lambda}\)。因此，交 \(N\cap M_{\lambda}\) 就是 N 的 \(\lambda\) 型同型分量。

对每个 \(\lambda \in \mathcal{S}\)，模 \(\mathrm{M}_{\lambda} + \mathrm{N} / \mathrm{N}\) 同构于 \(\mathrm{M}_{\lambda} / (\mathrm{M}_{\lambda} \cap \mathrm{N})\)。因而它是 \(\lambda\) 型同型模，且包含于 \((\mathrm{M} / \mathrm{N})_{\lambda}\)。最后的断言于是由 a) 及 II, §1, No.~8, 第 208 页的注 1 得出。

推论. —— 设 A 是一个环，S 是单 A-模的类组成的集合。设 M 是一个半单 A-模，N 是 M 的一个子模。则我们有 \(N = \bigoplus_{\lambda \in \mathscr{S}} N \cap M_{\lambda}\) 与 \(M/N = \bigoplus_{\lambda \in \mathscr{S}} (M_{\lambda} + N)/N\)。

由于 N 与 M/N 都是半单的（VIII, 第 56 页，推论 3），本推论由命题 4, d) 得出。

半单 A-模 M 的支集是使 M 的 \(\lambda\) 型同型分量非零的单 A-模的类 \(\lambda\) 组成的集合。有限生成的半单 A-模的支集是有限的。

命题 5. —— 设 A 是一个环，S 是单 A-模的类组成的集合。设 M 与 N 是 A-模。

\begin{enumerate}
\def\labelenumi{\alph{enumi})}
\item
  设 \(f: \mathrm{M} \to \mathrm{N}\) 是一个同态。对每个 \(\lambda \in \mathcal{S}\)，\(f\) 诱导出一个同态 \(f_{\lambda}\)（从 \(\mathrm{M}_{\lambda}\) 到 \(\mathrm{N}_{\lambda}\)）；若 \(\mathrm{M}\) 是半单的且 \(f\) 是满射，则每个同态 \(f_{\lambda}\) 都是满射。
\item
  设 M 是半单的。映射 \(f \mapsto (f_{\lambda})_{\lambda \in \mathcal{S}}\) 是从 \(\mathrm{Hom}_{\mathrm{A}}(\mathrm{M},\mathrm{N})\) 到 \(\prod_{\lambda \in \mathcal{S}}\mathrm{Hom}_{\mathrm{A}}(\mathrm{M}_{\lambda},\mathrm{N}_{\lambda})\) 的群同构。当 M 等于 N 时，该映射是从 \(\mathrm{End}_{\mathrm{A}}(\mathrm{M})\) 到 \(\prod_{\lambda \in \mathcal{S}}\mathrm{End}_{\mathrm{A}}(\mathrm{M}_{\lambda})\) 的环同构。
\end{enumerate}

对每个 \(\lambda \in \mathcal{S}\)，子模 \(f(\mathrm{M}_{\lambda})\)（\(\mathbf{N}\) 的）同构于某个 \(\lambda\) 型同型模的商；因而它是 \(\lambda\) 型同型模，从而包含于 \(\mathrm{N}_{\lambda}\)。

设 \(\mathrm{M}\) 是半单的且 \(f\) 是满射。于是 \(f\) 诱导出从 \(\mathrm{M} / \mathrm{Ker}(f)\) 到 \(\mathrm{N}\) 的同构，它把 \((\mathrm{M}_{\lambda} + \mathrm{Ker}(f)) / \mathrm{Ker}(f)\) 映到 \(f(\mathrm{M}_{\lambda})\)。

由 VIII, 第 65 页的命题 4，我们有 \(N_{\lambda} = f(M_{\lambda})\)，这就完成了 a) 的证明。

b) 中所考虑的映射显然是群同态，且当 M 等于 N 时是环同态。设 \((f_{\lambda})_{\lambda \in \mathcal{S}}\) 是 \(\prod_{\lambda \in \mathcal{S}} \mathrm{Hom}(M_{\lambda}, N_{\lambda})\) 的一个元素。它在 b) 中映射下的逆像的唯一元素是由下式定义的同态 \(f : M \to N\)：

\[
f \Bigl (\sum_ {\lambda \in \mathcal {S}} x _ {\lambda} \Bigr) = \sum_ {\lambda \in \mathcal {S}} f _ {\lambda} (x _ {\lambda})
\]

其中 \((x_{\lambda})_{\lambda \in \mathcal{S}}\in \bigoplus_{\lambda}\mathrm{M}_{\lambda}\) 任意。

注. —— 设 A 与 B 是环。设 M 是一个 (A, B)-双模。由命题 5 可知，A-模 M 的诸同型分量都是 M 的子双模。特别地，当 M 是一个 A-模且 B 是 \(\operatorname{End}_{\mathrm{A}}(\mathrm{M})\) 的反环时就属于这种情形。

例. —— 我们来考虑环 A 交换的情形。把极大理想 m 映到 \(\operatorname{cl}(A/m)\) 的映射是从 A 的极大理想组成的集合到单 A-模的类组成的集合 S 的一个双射（VIII, 第 51 页）。逆双射把 \(\lambda\) 映到它的零化理想 \(m_{\lambda}\)。

设 M 是一个 A-模。对每个 \(\lambda\in S\)，M 的同型分量 \(M_{\lambda}\)（\(\lambda\) 型）由被 \(m_{\lambda}\) 零化的元素组成，并且我们可以把 \(M_{\lambda}\) 看作体 \(A/m_{\lambda}\) 上的一个向量空间。若 M 是半单的而 N 是另一个 A-模，则由命题 5 可以得到从 \(\operatorname{Hom}_{\mathrm{A}}(\mathrm{M},\mathrm{N})\) 到 \(\prod_{\lambda\in\mathcal{S}}\operatorname{Hom}_{\mathrm{A}/m_{\lambda}}(\mathrm{M}_{\lambda},\mathrm{N}_{\lambda})\) 的一个群同构。

推论 1. —— 设 M 是一个半单 A-模，N 是 M 的一个子模。下列性质等价：

\begin{enumerate}
\def\labelenumi{(\roman{enumi})}
\item
  在 M 中存在唯一的与 N 互补的子模。
\item
  我们有 \(\mathrm{Hom}_{\mathrm{A}}(\mathrm{M} / \mathrm{N},\mathrm{N}) = 0\)
\item
  存在一个子集 \(\Lambda\)（含于 \(\mathcal{S}\)），使得 \(N = \bigoplus_{\lambda \in \Lambda} M_{\lambda}\)。
\end{enumerate}

选取 M 中与 N 互补的一个子模 \(N^{\prime}\)（VIII, 第 56 页，推论 2）。若我们把 M 与 \(N^{\prime} \times N\) 等同起来，则 M 中与 N 互补的子模就是从 \(N^{\prime}\) 到 N 的 A-线性映射的图象。由于 \(N^{\prime}\) 同构于 M/N，我们便证明了性质 (i) 与 (ii) 等价。由命题 5, b)，群 \(\operatorname{Hom}_{\mathrm{A}}(\mathrm{N}^{\prime}, \mathrm{N})\) 同构于群 \(\prod_{\lambda \in \mathcal{S}} \operatorname{Hom}_{\mathrm{A}}(\mathrm{N}_{\lambda}^{\prime}, \mathrm{N}_{\lambda})\)。它为零当且仅当对每个 \(\lambda \in S\) 有 \(N_{\lambda} = 0\) 或 \(N_{\lambda}^{\prime} = 0\)（VIII, 第 61 页，注），亦即 \(N_{\lambda} = 0\) 或 \(N_{\lambda} = M_{\lambda}\)。这就证明了性质 (ii) 与 (iii) 等价。

推论 2. —— 设 M 是一个 A-模。下列两个条件等价：

\begin{enumerate}
\def\labelenumi{(\roman{enumi})}
\item
  M 的每个子模都有唯一的补子模。
\item
  模 M 是一族两两不同构的单模 \((\mathrm{S}_{i})_{i\in\mathrm{I}}\) 的直和。
\end{enumerate}

设 M 满足这些条件。于是，对 M 的每个子模 N，存在 I 的唯一子集 J 使我们有 \(N = \bigoplus_{j \in J} S_j\)，且 M 的每个单子模都等于某个 \(S_i\)。

条件 (i) 与 (ii) 都蕴含 M 是半单的。

设条件 (i) 成立。设 \(\lambda \in \mathcal{S}\)。由推论 1 中 (i) 与 (iii) 的等价性，\(\mathrm{M}_{\lambda}\) 的每个子模或为零或等于 \(\mathrm{M}_{\lambda}\)。因此 \(\mathrm{M}_{\lambda}\) 或为零或为单模，而 (ii) 由 \(\mathrm{M}\) 是族 \((\mathrm{M}_{\lambda})_{\lambda \in \mathcal{S}}\) 的直和这一事实得出。

反之，若条件 (ii) 成立，则 \(M_{\lambda}\) 对每个 \(\lambda \in S\) 均为零或单模。若 N 是 M 的一个子模，则 \(N \cap M_{\lambda}\) 等于 0 或 \(M_{\lambda}\)（对每个 \(\lambda \in S\)）。由于我们有 \(N = \bigoplus_{\lambda \in S} (N \cap M_{\lambda})\)（VIII, 第 66 页，推论），子模 N 具有推论 1 的性质 (iii)，从而在 M 中有唯一的补子模。这就证明了 (ii) 蕴含 (i)，同时也证明了本推论的各条断言。

设 M 是一个 A-模，S 是一个单 A-模。用 D 表示体 \(\operatorname{End}_{\mathrm{A}}(\mathrm{S})\) 的反环，并把 S 视为 (A,D)-双模。于是 \(\operatorname{Hom}_{\mathrm{A}}(\mathrm{S},\mathrm{M})\) 是 D 上的左向量空间，而 \(\operatorname{Hom}_{\mathrm{A}}(\mathrm{M},\mathrm{S})\) 是 D 上的右向量空间。左 D-向量空间 \(\operatorname{Hom}_{\mathrm{A}}(\mathrm{S},\mathrm{M})\) 的对偶是 D 上的右向量空间（II, §2, No.~3, 第 232 页，定义 2）。对每个 \(u \in \operatorname{Hom}_{\mathrm{A}}(\mathrm{M},\mathrm{S})\)，映射 \(h(u): v \mapsto u \circ v\)（从 \(\operatorname{Hom}_{\mathrm{A}}(\mathrm{S},\mathrm{M})\) 到 \(\operatorname{Hom}_{\mathrm{A}}(\mathrm{S},\mathrm{S}) = \mathrm{D}\)）是左 D-向量空间 \(\operatorname{Hom}_{\mathrm{A}}(\mathrm{S},\mathrm{M})\) 上的线性型。

命题 6. —— 沿用上面的记号，并设 A-模 M 是半单的。映射 \(u \mapsto h(u)\)（从右 D-向量空间 \(\mathrm{Hom}_{\mathrm{A}}(\mathrm{M},\mathrm{S})\) 到左 D-向量空间 \(\mathrm{Hom}_{\mathrm{A}}(\mathrm{S},\mathrm{M})\) 的对偶）是 D-线性的双射。

设 \(u \in \mathrm{Hom}_{\mathrm{A}}(\mathrm{M}, \mathrm{S})\)，\(v \in \mathrm{Hom}_{\mathrm{A}}(\mathrm{S}, \mathrm{M})\)，\(d \in \mathrm{D}\)。我们有

\[
h (u d) (v) = h (d \circ u) (v) = d \circ u \circ v = d \circ (h (u) (v)) = h (u) (v) d.
\]

这就证明映射 h 是 D-线性的。它就是由 \(u \mapsto \operatorname{Hom}(1_{\mathrm{S}}, u)\) 给出的从 \(\operatorname{Hom}_{\mathrm{A}}(\mathrm{M}, \mathrm{S})\) 到 \(\operatorname{Hom}_{\mathrm{D}}(\operatorname{Hom}_{\mathrm{A}}(\mathrm{S}, \mathrm{M}), \operatorname{Hom}_{\mathrm{A}}(\mathrm{S}, \mathrm{S}))\) 的映射。要证明它是双射，由 VIII, 第 66 页的命题 5, b)，只需处理 M 是 S 型同型模的情形；这时可以应用 VIII, 第 65 页的推论。

\subsection{半单模的描述}

在本节余下部分，A 是一个环，S 是单 A-模的类组成的集合。对每个 \(\lambda \in S\)，我们选取一个单模 \(S_{\lambda}\)，其类为 \(\lambda\)（例如 \(S_{\lambda} = \lambda\)）；把 \(S_{\lambda}\) 的自同态体之反环记作 \(D_{\lambda}\)。我们把 \(S_{\lambda}\) 视为 \((A, D_{\lambda})\)-双模。

设 \(M\) 是一个 A-模。对每个 \(\lambda \in \mathcal{S}\)，\(\mathrm{Hom}_{\mathrm{A}}(\mathrm{S}_{\lambda}, M)\) 是体 \(D_{\lambda}\) 上的左向量空间。由 VIII, 第 59 页及 II, §1, No.~6, 第 202 页的命题 6，存在唯一的 A-线性映射，称为典范映射，

\[
\alpha_ {\mathrm{M}}: \bigoplus_ {\lambda \in \mathcal {S}} \left(\mathrm{S} _ {\lambda} \otimes_ {\mathrm{D} _ {\lambda}} \operatorname{Hom} _ {\mathrm{A}} (\mathrm{S} _ {\lambda}, \mathrm{M})\right) \to \mathrm{M}
\]

它满足关系式

\[
\alpha_ {\mathrm{M}} (s \otimes f) = f (s)\tag{10}
\]

其中 \(\lambda \in \mathcal{S}\)，\(s \in S_{\lambda}\)，\(f \in \mathrm{Hom}_{\mathrm{A}}(\mathrm{S}_{\lambda}, \mathrm{M})\)。若给 \(\bigoplus_{\lambda \in \mathcal{S}} (\mathrm{S}_{\lambda} \otimes_{\mathrm{D}_{\lambda}} \mathrm{Hom}_{\mathrm{A}}(\mathrm{S}_{\lambda}, \mathrm{M}))\) 与 \(\mathrm{M}\) 赋以它们自然的 \(\mathrm{End}_{\mathrm{A}}(\mathrm{M})\)-模结构，则映射 \(\alpha_{\mathrm{M}}\) 是 \(\mathrm{End}_{\mathrm{A}}(\mathrm{M})\)-线性的。

命题 7. —— 设 M 是一个 A-模。典范映射 \(\alpha_{\mathrm{M}}\) 是单射。对每个 \(\lambda \in \mathcal{S}\)，在 \(\alpha_{\mathrm{M}}\) 下，\(\mathrm{S}_{\lambda} \otimes_{\mathrm{D}_{\lambda}} \mathrm{Hom}_{\mathrm{A}}(\mathrm{S}_{\lambda}, \mathrm{M})\) 的像是 M 的 \(\lambda\) 型同型分量。\(\alpha_{\mathrm{M}}\) 的像是 M 的底座。A-模 M 是半单的当且仅当映射 \(\alpha_{\mathrm{M}}\) 是双射。

设 \(\lambda \in \mathcal{S}\)。把 M 的 \(\lambda\) 型同型分量记作 \(\mathrm{M}_{\lambda}\)。从 \(\mathrm{S}_{\lambda}\) 到 M 的每个 A-线性映射都取值于 \(\mathrm{M}_{\lambda}\)（VIII, 第 66 页，命题 5）。于是，由 VIII, 第 62 页的命题 3, a)，映射 \(\alpha_{\mathrm{M}}\) 诱导出从 \(\mathrm{S}_{\lambda} \otimes_{\mathrm{D}_{\lambda}} \mathrm{Hom}_{\mathrm{A}}(\mathrm{S}_{\lambda}, \mathrm{M})\) 到 \(\mathrm{M}_{\lambda}\) 的双射。由于 M 的底座是族 \((\mathrm{M}_{\lambda})_{\lambda \in \mathcal{S}}\) 的直和，且模 M 是半单的当且仅当它等于自身的底座，命题得证。

定义 5. —— 设 M 是一个半单 A-模。M 的（关于族 \((\mathrm{S}_{\lambda})_{\lambda\in\mathcal{S}}\) 的）描述是一个偶 \(((\mathrm{V}_{\lambda})_{\lambda\in\mathcal{S}},\alpha)\)，其中 \(V_{\lambda}\) 是体 \(D_{\lambda}\) 上的左向量空间（对每个 \(\lambda\in S\)），而 \(\alpha:\bigoplus_{\lambda\in\mathcal{S}}(\mathrm{S}_{\lambda}\otimes_{\mathrm{D}_{\lambda}}\mathrm{V}_{\lambda})\to\mathrm{M}\) 是 A-模的同构。

由命题 7，每个半单 A-模 M 都有一个典范描述：它是偶 \(\left(\left(\mathrm{Hom}_{\mathrm{A}}\left(\mathrm{S}_{\lambda},\mathrm{M}\right)\right)_{\lambda \in \mathcal{S}},\alpha_{\mathrm{M}}\right)\)，其中 \(\alpha_{\mathrm{M}}\) 是由公式 (10) 定义的 A-线性映射。

命题 8. —— 设 M 是一个半单 A-模，\(((\mathrm{V}_{\lambda})_{\lambda \in \mathcal{S}}, \alpha)\) 是 M 的一个描述。

\begin{enumerate}
\def\labelenumi{\alph{enumi})}
\item
  对每个 \(\lambda \in \mathcal{S}\)，映射 \(\alpha\) 诱导出从 A-模 \(S_{\lambda} \otimes_{D_{\lambda}} V_{\lambda}\) 到 M 的 \(\lambda\) 型同型分量的同构。
\item
  对每个 \(\lambda \in \mathcal{S}\)，映射 \(\beta_{\lambda} : \mathrm{V}_{\lambda} \to \mathrm{Hom}_{\mathrm{A}}(\mathrm{S}_{\lambda}, \mathrm{M})\)（由 \(\beta_{\lambda}(v)(s) = \alpha(s \otimes v)\) 定义）是 \(\mathrm{D}_{\lambda}\)-线性的双射。
\item
  设 \(\mathbf{N}\) 是 \(\mathbf{M}\) 的一个子模。存在唯一的族 \((\mathrm{W}_{\lambda})_{\lambda \in \mathcal{S}}\) 具有下列性质：\(\mathrm{W}_{\lambda}\) 是一个 \(\mathrm{D}_{\lambda}\)-线性子空间，含于 \(\mathrm{V}_{\lambda}\)（对每个 \(\lambda \in \mathcal{S}\)），且 \(\mathbf{N}\) 是 \(\alpha\) 作用于模 \(\bigoplus_{\lambda \in \mathcal{S}} (\mathrm{S}_{\lambda} \otimes_{\mathrm{D}_{\lambda}} \mathrm{W}_{\lambda})\) 所得的像，这里该模与其在模 \(\bigoplus_{\lambda \in \mathcal{S}} (\mathrm{S}_{\lambda} \otimes_{\mathrm{D}_{\lambda}} \mathrm{V}_{\lambda})\) 中的典范像等同。对每个 \(\lambda \in \mathcal{S}\)，\(\mathrm{W}_{\lambda}\) 是这样的元素 \(v \in \mathrm{V}_{\lambda}\) 组成的集合：\(\alpha(s \otimes v)\) 属于 \(\mathbf{N}\)（对每个 \(s \in \mathrm{S}_{\lambda}\)）。
\end{enumerate}

A-模 \(S_{\lambda} \otimes_{D_{\lambda}} V_{\lambda}\) 都是 \(\lambda\) 型同型模（对每个 \(\lambda \in S\)）。于是断言 a) 由 \(\alpha\) 是同构以及 M 是族 \((M_{\lambda})_{\lambda \in \mathcal{S}}\) 的直和（VIII, 第 65 页，命题 4, a)）这两个事实得出。

设 \(\lambda \in \mathcal{S}\)。把 \(\alpha\) 在 \(S_{\lambda} \otimes_{D_{\lambda}} V_{\lambda}\) 上的限制记作 \(\alpha_{\lambda}: S_{\lambda} \otimes_{D_{\lambda}} V_{\lambda} \to M\)。把 No.~3 的记号应用于 \((A, D_{\lambda})\)-双模 \(S_{\lambda}\)，我们有

\[
\beta_ {\lambda} = \gamma (\alpha_ {\lambda}) = \mathcal {H} (\alpha_ {\lambda}) \circ \beta_ {\mathrm{V} _ {\lambda}},
\]

其中最后一式由 VIII, 第 60 页得出。于是，\(\beta_{\lambda}\) 是 \(D_{\lambda}\)-线性同态 \(\mathcal{H}(\alpha_{\lambda})\) 与 \(\beta_{V_{\lambda}}\) 的复合。于是断言 b) 由 VIII, 第 62 页的命题 3, b) 得出。

设 \(N\) 是 \(M\) 的一个子模。我们有 \(N = \bigoplus_{\lambda \in \mathcal{S}} (N \cap M_{\lambda})\)（VIII, 第 66 页，推论），故 c) 由 VIII, 第 62 页的定理 2 得出。

推论. —— 设 M 是一个半单 A-模。对 M 的每个子模 N 及 S 的每个元素 \(\lambda\)，把 \(\operatorname{Hom}_{\mathrm{A}}(\mathrm{S}_{\lambda},\mathrm{N})\) 与一个 \(D_{\lambda}\)-线性子空间等同起来，该子空间由 \(\operatorname{Hom}_{\mathrm{A}}(\mathrm{S}_{\lambda},\mathrm{M})\) 中像含于 N 的映射组成。

\begin{enumerate}
\def\labelenumi{\alph{enumi})}
\item
  映射 \(\mathrm{N}\mapsto(\mathrm{Hom}_{\mathrm{A}}(\mathrm{S}_{\lambda},\mathrm{N}))_{\lambda\in\mathcal{S}}\) 是从 M 的 A-子模组成的集合到满足下列条件的族 \((\mathrm{W}_{\lambda})_{\lambda\in\mathcal{S}}\) 组成的集合的双射：对每个 \(\lambda\in S\)，\(W_{\lambda}\) 是一个 \(D_{\lambda}\)-线性子空间，含于 \(\mathrm{Hom}_{\mathrm{A}}(\mathrm{S}_{\lambda},\mathrm{M})\)。
\item
  逆双射把族 \((\mathrm{W}_{\lambda})_{\lambda \in \mathcal{S}}\) 映到 M 的 A-子模 \(\sum_{\lambda \in \mathcal{S}}\sum_{w\in \mathrm{W}_{\lambda}}w(\mathrm{S}_{\lambda})\)。
\end{enumerate}

这是把命题 8, c) 应用于 M 的典范描述的一个重新表述。

命题 9. —— 设 M 与 M' 是半单 A-模，并分别考虑 M 与 M' 的描述 \(((\mathrm{V}_{\lambda})_{\lambda\in\mathcal{S}},\alpha)\) 与 \(((\mathrm{V}_{\lambda}^{\prime})_{\lambda\in\mathcal{S}},\alpha^{\prime})\)。对每个族 \(\boldsymbol{f}=(f_{\lambda})_{\lambda\in\mathcal{S}}\)（取自 \(\prod_{\lambda\in\mathcal{S}}\mathrm{Hom}_{\mathrm{D}_{\lambda}}(\mathrm{V}_{\lambda},\mathrm{V}_{\lambda}^{\prime})\)），存在唯一的 A-线性映射 \(\varphi(\boldsymbol{f})\in\mathrm{Hom}_{\mathrm{A}}(\mathrm{M},\mathrm{M}^{\prime})\) 使下面的图式交换：

\[
\begin{array}{c} \bigoplus_ {\lambda \in \mathscr {S}} (S _ {\lambda} \otimes_ {D _ {\lambda}} V _ {\lambda}) \xrightarrow {\alpha} M \\ \bigoplus_ {\lambda \in \mathscr {S}} (1 _ {S _ {\lambda}} \otimes f _ {\lambda}) \Bigg | _ {\downarrow} \qquad \qquad \qquad \qquad \qquad \qquad \Bigg | _ {\varphi (\boldsymbol {f})} \\ \bigoplus_ {\lambda \in \mathscr {S}} (S _ {\lambda} \otimes_ {D _ {\lambda}} V _ {\lambda} ^ {\prime}) \xrightarrow {\alpha^ {\prime}} M ^ {\prime}. \end{array}
\]

这样定义的映射 \(\varphi : \prod_{\lambda \in \mathcal{S}} \mathrm{Hom}_{\mathrm{D}_{\lambda}}(\mathrm{V}_{\lambda}, \mathrm{V}_{\lambda}') \to \mathrm{Hom}_{\mathrm{A}}(\mathrm{M}, \mathrm{M}')\) 是群同构。当我们有 \(\mathrm{M} = \mathrm{M}'\)、\(\mathrm{V}_{\lambda} = \mathrm{V}_{\lambda}'\)（对每个 \(\lambda \in \mathcal{S}\)）、且 \(\alpha = \alpha'\) 时，映射 \(\varphi\) 是从 \(\prod_{\lambda \in \mathcal{S}} \mathrm{End}_{\mathrm{D}_{\lambda}}(\mathrm{V}_{\lambda})\) 到 \(\mathrm{End}_{\mathrm{A}}(\mathrm{M})\) 的环同构。

鉴于命题 8, a) 给出的 M 与 M' 的同型分量的描述，这由 VIII, 第 64 页的定理 3 及 VIII, 第 66 页的命题 5, b) 得出。

推论. —— 设 M 是一个半单 A-模，M' 是一个 A-模。映射 \(u \mapsto (\operatorname{Hom}(1_{\mathrm{S}_{\lambda}}, u))_{\lambda \in \mathcal{S}}\) 从 \(\operatorname{Hom}_{\mathrm{A}}(\mathrm{M}, \mathrm{M}')\) 到

\[
\prod_ {\lambda \in \mathcal {S}} \mathrm{Hom} _ {\mathrm{D} _ {\lambda}} (\mathrm{Hom} _ {\mathrm{A}} (\mathrm{S} _ {\lambda}, \mathrm{M}), \mathrm{Hom} _ {\mathrm{A}} (\mathrm{S} _ {\lambda}, \mathrm{M} ^ {\prime}))
\]

是群同构。当 \(\mathrm{M}'\) 等于 \(\mathrm{M}\) 时，它是从环 \(\operatorname{End}_{\mathrm{A}}(\mathrm{M})\) 到环 \(\prod_{\lambda \in \mathcal{S}} \operatorname{End}_{\mathrm{D}_{\lambda}}(\operatorname{Hom}_{\mathrm{A}}(\mathrm{S}_{\lambda}, \mathrm{M}))\) 的同构。

这是把命题 9 应用于 M 的典范描述与 M' 的底座的典范描述的一个重新表述。

\subsection{半单模中的重数与长度}

命题 10. —— 设 M 是一个半单 A-模。设 \((\mathrm{M}_{i})_{i\in\mathrm{I}}\) 是以 M 为直和的单子模族。下列性质等价：

\begin{enumerate}
\def\labelenumi{(\roman{enumi})}
\item
  M 的长度有限。
\item
  M 是阿廷的。
\item
  M 是诺特的。
\item
  M 是有限生成的。
\item
  I 是有限的。
\end{enumerate}

若 M 具有这些性质，则 M 的长度等于 I 的基数。

若集合 I 有限，则 M 具有性质 (i)、(ii)、(iii)、(iv)。设集合 I 无限。由 VIII, 第 2 页的例 2，模 M 既非阿廷的亦非诺特的；由于每个有限长度的模都是阿廷的且诺特的（VIII, 第 2 页，命题 1），M 也不是有限长度的。最后，M 的每个元素都属于有限个子模 \(M_{i}\) 之和，故 M 不是有限生成的。这就证明了性质 (i) 至 (v) 的等价性。若这些性质成立，则我们有 \(\text{long}(M) = \sum_{i \in I} \text{long}(M_{i}) = \text{Card}(I)\)（II, §1, No.~10, 第 213 页，推论 5）。

命题 11. —— 设 M 是一个半单 A-模，它是单子模族 \((\mathrm{M}_{i})_{i\in\mathrm{I}}\) 的直和。对每个 \(\lambda\in S\)，用 \(\mathrm{I}(\lambda)\) 表示这样的指标 \(i\in I\) 组成的集合：\(M_{i}\) 属于类 \(\lambda\)。\(\mathrm{I}(\lambda)\) 的基数等于左 \(D_{\lambda}\)-向量空间 \(\mathrm{Hom}_{\mathrm{A}}(\mathrm{S}_{\lambda},\mathrm{M})\) 的维数。

A 的 \(\lambda\) 型同型分量同构于 \(\mathrm{S}_{\lambda}^{(\mathrm{I}(\lambda))}\)（VIII, 第 65 页，命题 4, b)）。\(D_{\lambda}\)-向量空间 \(\mathrm{Hom}_{\mathrm{A}}(\mathrm{S}_{\lambda},\mathrm{M})\) 可与 \(\mathrm{Hom}_{\mathrm{A}}(\mathrm{S}_{\lambda},\mathrm{M}_{\lambda})\) 等同，故同构于 \(\mathrm{D}_{\lambda}^{(\mathrm{I}(\lambda))}\)（No.~2）。这就证明了本命题。

每个单模都是原始模（VIII, 第 45 页），故每个半单模都是半原始模。设 M 是一个半单 A-模，\(\lambda \in \mathcal{S}\)。\(\lambda\) 在 M 中的重数是 VIII, 第 34 页定义的 \([\mathrm{M}:\lambda]\)，称为 \(\lambda\) 在 M 中的原始重数。命题 11 可以转述为等式

\[
[ \mathrm{M}: \lambda ] = \dim_ {\mathrm{D} _ {\lambda}} (\operatorname{Hom} _ {\mathrm{A}} (\mathrm{S} _ {\lambda}, \mathrm{M})).\tag{11}
\]

更一般地，若 \(((\mathrm{V}_{\lambda})_{\lambda \in \mathcal{S}},\alpha)\) 是 M 的一个描述，则 \([\mathrm{M}:\lambda ]\) 等于 \(\dim_{\mathrm{D}_{\lambda}}(\mathrm{V}_{\lambda})\)。由 VIII, 第 68 页的命题 6，我们还有

\[
[ \mathrm{M}: \lambda ] = \dim_ {\mathrm{D} _ {\lambda}} (\operatorname{Hom} _ {\mathrm{A}} (\mathrm{M}, \mathrm{S} _ {\lambda}))\tag{12}
\]

其中重数 \([M : \lambda]\) 有限。半单 A-模 M 与 \(M'\) 同构当且仅当我们有 \([M : \lambda] = [M' : \lambda]\)（对每个 \(\lambda \in S\)）。

设 M 是一个半单 A-模。存在一个具有下列性质的基数 I：对 M 分解为单模直和的每个分解 \(M = \bigoplus_{i \in I} M_i\)，I 的基数都等于 I（VIII, 第 34 页，推论 2）。这个基数称为半单 A-模 M 的长度，记作 \(\operatorname{long}_{\mathrm{A}}(\mathrm{M})\) 或 \(\operatorname{long}(\mathrm{M})\)。当 M 的长度有限时，由命题 10，这个定义与 II, §1, No.~10, 第 212 页中的定义相容。

单 A-模是长度为 1 的半单 A-模，并且我们有公式

\[
\operatorname{long} _ {\mathrm{A}} \left(\bigoplus_ {j \in \mathrm{J}} \mathrm{M} _ {j}\right) = \sum_ {j \in \mathrm{J}} \operatorname{long} _ {\mathrm{A}} (\mathrm{M} _ {j})\tag{13}
\]

其中 \((\mathrm{M}_j)_{j\in \mathrm{J}}\) 是任意一族半单 A-模。由命题 11，我们有

\[
\mathrm{long} _ {\mathrm{A}} (\mathrm{M}) = \sum_ {\lambda \in \mathcal {S}} \dim_ {\mathrm{D} _ {\lambda}} \mathrm{Hom} _ {\mathrm{A}} (\mathrm{S} _ {\lambda}, \mathrm{M}).\tag{14}
\]

把这个公式应用于 \(M_{\lambda}\)，我们得到 \([M : \lambda] = \text{long}_{A}(M_{\lambda})\)（对每个 \(\lambda \in S\)）。

当 A 是一个体时，单 A-模是维数为 1 的向量空间；这时每个 A-模都是半单的（II, §7, No.~1, 第 292 页，定理 1），而它的长度就是它作为 A 上向量空间的维数（II, §7, No.~2, 第 293 页）。

\BourbakiExercises{习题}

\begin{enumerate}
\def\labelenumi{\arabic{enumi})}
\tightlist
\item
  \begin{enumerate}
  \def\labelenumii{\alph{enumii})}
  \tightlist
  \item
    设 M 是一个 A-模，x 是 M 中零化子为有限族极大理想之交的元素。证明模 Ax 是半单的（把它嵌入到单模的一个有限乘积中）。
  \end{enumerate}
\end{enumerate}

\begin{enumerate}
\def\labelenumi{\alph{enumi})}
\setcounter{enumi}{1}
\tightlist
\item
  由此推出：一个模 M 是半单的当且仅当每个非零元素的零化子都是有限族极大理想之交。
\end{enumerate}

\begin{enumerate}
\def\labelenumi{\arabic{enumi})}
\setcounter{enumi}{1}
\tightlist
\item
  设 M 是一个模。证明下列性质等价：
\end{enumerate}

\(\alpha)\) M 是半单的且长度有限。

\(\beta)\) M 是诺特的，且 M 的每个极大子模都有补。

\(\gamma)\) M 是阿廷的，且 M 的每个单子模都有补。

\begin{enumerate}
\def\labelenumi{\arabic{enumi})}
\setcounter{enumi}{2}
\item
  设 \(\mathrm{K}\) 是一个体，I 是无限集，A 是环 \(\mathrm{K}^{\mathrm{I}}\)。证明 A-模 \(\mathrm{A}_s\) 不是半单的，尽管每个单生成理想都有补（确定 \(\mathrm{A}_s\) 的所有单子模）。
\item
  设 M 是一个模，它是子模族 \((\mathrm{M}_{\alpha})_{\alpha\in\mathrm{I}}\) 的直和。证明子模 \(M_{\alpha}\) 在 M 的所有自同构下稳定当且仅当每个同态 \(M_{\alpha}\to M_{\beta}\)（其中 \(\alpha\neq\beta\)）都为零。
\item
  设 M 是一个有限长度的模，\(M = \oplus_{i \in I} M_i\) 是 M 分解为不可分解子模直和的一个分解。设 \(J \subset I\) 是 i 与 j 之间关系“模 \(M_i\) 与 \(M_j\) 同构”的一个等价类；令 \(M_J = \oplus_{i \in J} M_i\)。举一个例子表明 \(M_J\) 不一定在 M 的每个自同构下稳定（取 M 为有限长度的 Z-模并利用习题 4）。
\item
  若模 M 的每个非零商都含一个单子模，则称 M 为半阿廷模。每个阿廷模以及体上的每个向量空间都是半阿廷的。
\end{enumerate}

\begin{enumerate}
\def\labelenumi{\alph{enumi})}
\item
  设 \(0 \rightarrow M' \rightarrow M \rightarrow M'' \rightarrow 0\) 是一个正合列。则 M 是半阿廷的当且仅当 \(M'\) 与 \(M''\) 亦然（若 M 是半阿廷的，考虑由满足 \(N \cap M' = 0\) 的 M 的子模 N 组成的集合中的一个极大元素）。
\item
  一族半阿廷模的直和是半阿廷模。 c) 设 K 是一个体，A 是乘积环 \(K^{N}\)。证明 A-模 \(A_{s}\) 不是半阿廷的，尽管它是单 A-模的一个乘积（若 s 是 \(A_{s}\) 的底座，证明 A-模 \(A_{s}/s\) 不含任何单模）。
\item
  设 A 是左诺特环，M 是半阿廷 A-模。证明 M 是它的阿廷子模之和（设 S 是 M 的阿廷子模之和；注意，若 N 是 M 的在 M/S 中有单像的有限生成子模，则 \(N \cap S\) 是阿廷的）。特别地，若 M 是有限生成的，则它是阿廷的。
\item
  设 D 是一个体，V 是无限维 D-向量空间，A 是由 D 的中心与有限秩自同态生成的 \(\operatorname{End}_{\mathrm{D}}(\mathrm{V})\) 的子环。证明 \(A_{s}\) 是半阿廷的（利用 a) 与 b)），但不是它的阿廷子模之和。
\end{enumerate}

\begin{enumerate}
\def\labelenumi{\arabic{enumi})}
\setcounter{enumi}{6}
\tightlist
\item
  设 A 是一个环，M 是一个 A-模。若 M 的子模 N 对 M 的每个非零子模 P 都满足 \(N \cap P \neq 0\)，则称 N 为本质子模。
\end{enumerate}

\begin{enumerate}
\def\labelenumi{\alph{enumi})}
\item
  证明 M 的每个子模 N 都是某个本质子模的直因子（考虑一个在与 N 之交为零的诸子模中极大的子模 P）。
\item
  证明 M 的底座是 M 的诸本质子模之交（由 a) 推出这个交是半单的）。特别地，M 是半单的当且仅当它不含除自身以外的本质子模。
\end{enumerate}

¶ 8) 设 A 是左诺特环。

\begin{enumerate}
\def\labelenumi{\alph{enumi})}
\item
  设 x 是 A 的非零元素。证明 x 的左零化子 Ann x 不是本质的（在不具有此性质的元素中，取一个使 Ann x 极大的 x，并考虑 \(Ax \cap Ann x\)）。
\item
  设 x 是 A 的右可约元。证明理想 Ax 是本质的（若 a 是 A 的理想且满足 \(a \cap Ax = 0\)，考虑理想 \(a + ax + ax^{2} + \ldots\)）。
\item
  由 a) 与 b) 推出，A 的每个右可约元都是可约元。
\item
  设 a 是 A 的一个含非幂零元素的本质左理想。证明 a 含一个可约元（设 \((x_{1},\ldots,x_{n})\) 是 a 中满足 Ann \(x_{i}^{2}=\) Ann \(x_{i}\) 且 \(x_{i+1}\in\) Ann \(x_{i}\cap\cdots\cap\) Ann \(x_{1}\)（对每个 \(i\geqslant1\)）的非零元素序列；注意 a 含诸 \(Ax_{i}\) 的直和。由此推出存在这样的序列使得 \(\bigcap_{i}\) Ann \(x_{i}=0\)；证明元素 \(\sum x_{i}^{2}\) 这时是右可约的，并利用 c)）。
\end{enumerate}

\section{交换}

\subsection{模的交换子与双交换子}

设 E 是一个环，B 是 E 的一个子集。B 在 E 中的交换子（或中心化子）是 E 中由与 B 的每个元素交换的元素组成的子环 \(B'\)。把 \(B''\)（即 \(B'\) 的交换子）称为 B 的双交换子（或双中心化子）。我们有 \(B \subset B''\)，且 \(B'\) 与其双交换子重合（III, §1, No.~2, 第 429 页）。E 的子环 B 的中心是 \(B \cap B'\)；\(B'\) 与 \(B''\) 的公共中心是 \(B' \cap B''\)。若 B 是 E 的交换子环，则我们有 \(B \subset B'\)，且 \(B''\) 是 \(B'\) 的中心（III, §1, No.~2, 第 430 页）。

设 A 是一个环，M 是左（相应地，右）A-模。我们把这些定义应用于下述情形：E 是 M 的加群的自同态环，B 是 M 的位似环 \(A_{M}\)。M 的交换子是指 \(A_{M}^{\prime}\)，即 \(A_{M}\) 在 E 中的交换子；它是 A-模 M 的自同态环。M 的双交换子是指 \(A_{M}^{\prime\prime}\)，即 \(A_{M}\) 在 E 中的双交换子；它是 M 的反模的自同态环（VIII, 第 8 页，定义 3）。我们有 \(A_{M} \subset A_{M}^{\prime\prime}\)，环 \(A_{M} \cap A_{M}^{\prime}\) 是 \(A_{M}\) 的中心，而 \(A_{M}^{\prime} \cap A_{M}^{\prime\prime}\) 是 \(A_{M}^{\prime}\) 与 \(A_{M}^{\prime\prime}\) 的公共中心。

定义 1. —— 若我们有 \(\mathrm{A}_{\mathrm{M}} = \mathrm{A}_{\mathrm{M}}^{\prime \prime}\)，则称 A-模 M 是平衡的。

若 A-模 M 是平衡的，则环 \(A_{M}\) 与 \(A_{M}^{\prime}\) 有相同的中心 \(A_{M} \cap A_{M}^{\prime}\)。模 M 是平衡 A-模当且仅当它是平衡 \(A_{M}\)-模。A-模 M 的反模是忠实的且平衡的。

当环 A 交换时，M 的双交换子 \(A_{M}^{\prime\prime}\) 是 \(\mathrm{A}_{\mathrm{M}}^{\prime}=\mathrm{End}_{\mathrm{A}}(\mathrm{M})\) 的中心；M 是平衡的意思是 \(\mathrm{End}_{\mathrm{A}}(\mathrm{M})\) 的中心化为位似。

对 A 的每个元素 \(a\)，用 \(\delta_{a}\) 表示从 A 到 A 的右位似 \(x\mapsto xa\)，用 \(\gamma_{a}\) 表示左位似 \(x\mapsto ax\)（I, §8, No.~1, 第 97 页）。映射 \(a\mapsto \delta_{a}\) 是从 \(\mathrm{A}^{\circ}\) 到左 A-模 \(\mathrm{A}_s\) 的交换子的环同构（参见 II, §10, No.~7, 第 349 页应用于 \(\mathrm{I} = \{1\}\)）。映射 \(a\mapsto \gamma_{a}\) 是从 A 到右 A-模 \(\mathrm{A}_d\) 的交换子的环同构（同前所引）。若经此映射把 A 与 \(\mathrm{A}_d\) 的交换子等同起来，则 \(\mathrm{A}_d\) 的反模与 \(\mathrm{A}_s\) 等同；因此，A-模 \(\mathrm{A}_d\) 是平衡的。同样，A-模 \(\mathrm{A}_s\) 也是平衡的。

设 n 是整数 \(\geqslant 1\)。我们把 \(A^{n}\) 视为左 \(\mathbf{M}_{n}(\mathrm{A})\)-模（同前所引）。把 \(m \in \mathbf{M}_{n}(\mathrm{A})\) 映到 \(x \mapsto mx\)（A-模 \(A_{d}^{n}\) 的自同态）的映射是从 \(\mathbf{M}_{n}(\mathrm{A})\) 到右 A-模 \(A_{d}^{n}\) 的交换子的环同构（同前所引）。

命题 1. —— 设 \((\mathrm{A}_{i})_{i\in\mathrm{I}}\) 是一族环，且对每个 \(i\in I\)，设 \(M_{i}\) 是一个 \(A_{i}\)-模。令 \(A=\prod A_{i}, M=\prod M_{i}\)，\(N=\bigoplus M_{i}\)。给 M 赋以作用律为 \(((a_{i}),(x_{i}))\mapsto(a_{i}x_{i})\) 的 A-模结构。集合 N 是 M 的 A-子模。

\begin{enumerate}
\def\labelenumi{\alph{enumi})}
\item
  映射 \((u_i) \mapsto \prod u_i\)（从 \(\prod \mathrm{End}_{\mathbf{Z}}(\mathrm{M}_i)\) 到 \(\mathrm{End}_{\mathbf{Z}}(\mathrm{M})\)，II, §1, No.~5, 第 200 页）分别限制为从 \(\prod (\mathrm{A}_i)_{\mathrm{M}_i}\)、\(\prod (\mathrm{A}_i)'_{\mathrm{M}_i}\) 与 \(\prod (\mathrm{A}_i)'''_{\mathrm{M}_i}\) 到 \(\mathrm{A}_{\mathrm{M}}\)、\(\mathrm{A}_{\mathrm{M}}'\) 与 \(\mathrm{A}_{\mathrm{M}}''\) 的环同构。
\item
  映射 \((u_i) \mapsto \bigoplus u_i\)（从 \(\prod \operatorname{End}_{\mathbf{Z}}(\mathrm{M}_i)\) 到 \(\operatorname{End}_{\mathbf{Z}}(\mathrm{N})\)，II, §1, No.~6, 第 203 页）分别限制为从 \(\prod (\mathrm{A}_i)_{\mathrm{M}_i}\)、\(\prod (\mathrm{A}_i)'_{\mathrm{M}_i}\) 与 \(\prod (\mathrm{A}_i)'''_{\mathrm{M}_i}\) 到 \(\mathrm{A_N}\)、\(\mathrm{A_N}'\) 与 \(\mathrm{A_N}''\) 的环同构。
\end{enumerate}

命题 1 中定义的环同构称为典范同构。经这些同构，乘积 \(\prod (\mathrm{A}_i)_{\mathrm{M}_i}\) 常与 \(\mathrm{A_M}\) 等同，\(\prod (\mathrm{A}_i)'_{\mathrm{M}_i}\) 与 \(\mathrm{A_M}'\) 等同，等等。

映射 \(\varphi : (u_i) \mapsto \prod u_i\)（从 \(\prod \operatorname{End}_{\mathbf{Z}}(\mathrm{M}_i)\) 到 \(\operatorname{End}_{\mathbf{Z}}(\mathrm{M})\)）是单的环同态。由 M 上 A-模结构的定义，我们有 \(\varphi(\prod (\mathrm{A}_i)_{\mathrm{M}_i}) = \mathrm{A}_{\mathrm{M}}\)。设 \(u \in \mathrm{A}_{\mathrm{M}}'\)。对每个 \(i \in \mathrm{I}\)，用 \(h_i\) 表示 A 中除指标 \(i\) 的分量为 0 外其余分量均为 1 的元素。若 \(x\) 是 M 的一个其指标 \(i\) 的分量为 0 的元素，则我们有 \(x = h_i x\)，从而 \(\operatorname{pr}_i(u(x)) = \operatorname{pr}_i(u(h_i x)) = \operatorname{pr}_i(h_i u(x)) = 0\)。因此，存在唯一的群同态 \(u_i: \mathrm{M}_i \mapsto \mathrm{M}_i\) 使 \(\operatorname{pr}_i(u(y)) = u_i(\operatorname{pr}_i(y))\) 对每个 \(y \in \mathrm{M}\) 成立。我们有 \(u = \prod u_i\)。由于映射 \(u\) 是 A-线性的，故映射 \(u_i\) 是 \(\mathrm{A}_i\)-线性的（对每个 \(i \in \mathrm{I}\)）。这就证明 \(\prod (\mathrm{A}_i)'_{\mathrm{M}_i}\) 在 \(\varphi\) 下的像包含 \(\mathrm{A}_{\mathrm{M}}'\)；反向包含是明显的。把这一结果应用于 M 的反模，便推出 \(\varphi\) 诱导出从 \(A_{M}^{\prime \prime}\) 到 \(\prod_i(A_i)_{M_i}^{\prime \prime}\) 的同构。这就证明了断言 a)。

b) 的证明与 a) 的证明经必要修改后相同。

命题 2. —— 设 A 是一个环，M 是一个 A-模。设 I 是一个集合。则 A-模 \(\mathrm{M}^{(I)}\) 的双交换子与 \(A_{M}^{\prime\prime}\)-模 \(\mathrm{M}^{(I)}\) 的位似环重合。

对 \(i \in I\)，把从 \(M^{(I)}\) 到 M 的投影同态记作 \(\pi_{i} : (x_{j})_{j \in I} \mapsto x_{i}\)，把相应的典范单射（I, §4, No.~9, 第 47 页）记作 \(\iota_{i} : M \to M^{(I)}\)。

设 \(u\) 是 \(\mathrm{End}_{\mathrm{A}}(\mathrm{M}^{(\mathrm{I})})\) 的一个元素。对所有 \(i,j \in \mathrm{I}\)，复合 \(u_{i,j} = \pi_j \circ u \circ \iota_i\) 属于交换子 \(\mathrm{A}_{\mathrm{M}}'\)（\(\mathrm{M}\) 的）。对每个元素 \(b\)（属于 \(\mathrm{A}_{\mathrm{M}}''\)）及每个 \((x_i) \in \mathrm{M}^{(\mathrm{I})}\)，我们有关系式

\[
b u \left(\left(x _ {i}\right) _ {i \in \mathrm{I}}\right) = b \left(\sum_ {i \in \mathrm{I}} u _ {i, j} \left(x _ {i}\right)\right) _ {j \in \mathrm{I}} = \left(\sum_ {i \in \mathrm{I}} u _ {i, j} \left(b x _ {i}\right)\right) _ {j \in \mathrm{I}} = u \left(b \left(x _ {i}\right) _ {i \in \mathrm{I}}\right).
\]

因此位似 \(b_{\mathrm{M}^{(\mathrm{I})}}\) 属于 A-模 \(\mathrm{M}^{(\mathrm{I})}\) 的双交换子。

反之，设 b 是 \(\mathrm{M}^{(\mathrm{I})}\) 的双交换子的一个元素。对 \(i,j\in I\)，令 \(b_{i,j}=\pi_{j}\circ b\circ\iota_{i}\)。取 \(i,j\in I\) 且 \(i\neq j\)。由于 \(\iota_{j}\circ\pi_{j}\) 属于 A-模 \(\mathrm{M}^{(\mathrm{I})}\) 的交换子，我们有

\[
b _ {i, j} = \pi_ {j} \circ \iota_ {j} \circ \pi_ {j} \circ b \circ \iota_ {i} = \pi_ {j} \circ b \circ \iota_ {j} \circ \pi_ {j} \circ \iota_ {i} = 0.
\]

同样，我们有

\[
b _ {j, j} = \pi_ {j} \circ b \circ \iota_ {j} = \pi_ {i} \circ \iota_ {i} \circ \pi_ {j} \circ b \circ \iota_ {j} = \pi_ {i} \circ b \circ \iota_ {i} \circ \pi_ {j} \circ \iota_ {j} = b _ {i, i}.
\]

此外，\(b_{i,i}\) 属于 \(A_{M}^{\prime\prime}\)。由此推出，b 与 \(A_{M}^{\prime\prime}\)-模 \(M^{(I)}\) 的一个位似重合。

\subsection{生成模}

设 \(\mathrm{A}\) 是一个环。

定义 2. —— 若每个 A-模 N 都由从 M 到 N 的 A-线性映射的像生成，则称 A-模 M 为生成模。

设 \(M\) 是一个左 A-模。我们把 \(M\) 的对偶记作 \(M^*\)，并把 \(M \times M^*\) 上的典范双线性型（II, §2, No.~3, 第 233 页）记作

\[
(x, x ^ {*}) \mapsto \langle x, x ^ {*} \rangle = x ^ {*} (x).
\]

用 \(\tau(\mathrm{M})\) 表示 A 中形如 \(\sum_{i=1}^{n}\langle x_i,x_i^*\rangle\) 的元素组成的集合，其中 \(x_1,\ldots ,x_n\) 是 M 的元素，\(x_1^*,\ldots ,x_n^*\) 是 \(\mathrm{M}^*\) 的元素。它是 A 的一个双边理想，称为 M 的迹理想。A-模 \(\mathrm{A}_s\) 的迹理想是 A。一族 A-模 \((\mathrm{M}_i)_{i\in \mathrm{I}}\) 的直和的迹理想是理想 \(\sum_{i\in \mathrm{I}}\tau(\mathrm{M}_i)\)。若 M 是投射 A-模，则由 II, §2, No.~6, 第 238 页的命题 12，我们有 \(\mathrm{M} = \tau(\mathrm{M})\mathrm{M}\)。

定理 1. —— 设 M 是一个左 A-模。下列性质等价：

\begin{enumerate}
\def\labelenumi{(\roman{enumi})}
\item
  A-模 M 是生成模。
\item
  对每个 A-模 N，存在一个集合 I 和一个从 \(M^{(I)}\) 到 N 的满 A-线性映射。
\item
  存在一个整数 \(n \geqslant 0\) 和一个从 \(\mathbf{M}^n\) 到 \(\mathrm{A}_s\) 的满 A-线性映射。
\item
  存在整数 \(n \geqslant 0\)，使 \(A_s\) 同构于 \(M^n\) 的一个直因子子模。
\item
  迹理想 \(\tau (\mathbf{M})\) 等于 A。
\item
  存在整数 \(n \geqslant 0\)、M 的元素 \(x_{1}, \ldots, x_{n}\) 及 M* 的元素 \(x_{1}^{*}, \ldots, x_{n}^{*}\)，满足 \(\sum_{i=1}^{n} \langle x_{i}, x_{i}^{*} \rangle = 1\)。
\item
  \(\Rightarrow\) (ii)：设 M 是生成模，N 是一个左 A-模。存在一族从 M 到 N 的 A-线性映射 \((u_{i})_{i\in\mathrm{I}}\)，使我们有 \(N=\sum_{i\in\mathrm{I}}u_{i}(\mathrm{M})\)。映射 \((x_{i})\mapsto\sum_{i\in\mathrm{I}}u_{i}(x_{i})\)（从 \(\mathrm{M}^{(\mathrm{I})}\) 到 N）是 A-线性的满射。
\item
  \(\Rightarrow\) (iii)：设性质 (ii) 成立，取 \(N = A_s\)。于是我们有一个集合 I 和一个满线性映射 \(u: M^{(I)} \to A_s\)。由于 \(A_s\) 由元素 1 生成，存在 I 的有限子集 J 使 \(u(M^{(J)}) = A_s\)，从而得到 (iii)。
\item
  \(\Rightarrow\) (iv)：这由 II, §1, No.~12, 第 218 页的命题 21 得出。
\item
  \(\Rightarrow\) (v)：设 \(n \geqslant 1\) 是使 \(\mathrm{A}_s\) 同构于 M 的一个直因子子模的整数。我们有 \(\mathrm{A} = \tau(\mathrm{A}_s) \subset \tau(\mathrm{M}^n) = \tau(\mathrm{M})\)，因而 \(\tau(\mathrm{M}) = \mathrm{A}\)。
\item
  \(\Rightarrow\) (vi)：这是明显的。
\item
  \(\Rightarrow\) (i)：设 \(n\) 是整数 \(\geqslant 0\)，\(x_{1},\ldots ,x_{n}\) 是 M 的元素，\(x_{1}^{*},\ldots ,x_{n}^{*}\) 是 M \(^*\) 的元素，满足 \(\sum_{i = 1}^{n}\langle x_i,x_i^*\rangle = 1\)。设 N 是一个左 A-模，\(y\) 是 N 的一个元素。从 M 到 N 的映射 \(u_{i}:x\mapsto \langle x,x_{i}^{*}\rangle y\) 是 A-线性的，且我们有 \(y = \sum_{i = 1}^{n}u_{i}(x_{i})\)。这就证明 M 是一个生成 A-模。
\end{enumerate}

推论. —— 生成 A-模是忠实的。

设 \(a \in A\) 满足 \(aM = 0\)。利用定理 1 的蕴含关系 (i) \(\Rightarrow\) (iv)，我们得到 \(aA_s = 0\)，因此 \(a = 0\)。推论得证（II, §1, No.~12, 第 219 页）。

例. —— 1) A-模 \(A_s\) 是生成模。

\begin{enumerate}
\def\labelenumi{\arabic{enumi})}
\setcounter{enumi}{1}
\item
  每个非零自由 A-模都是生成模。更一般地，有生成商的模本身是生成模。
\item
  设 M 是一个反模为有限生成的半单 A-模。则 M 是生成 \(A_{M}\)-模。事实上，由 VIII, 第 8 页的引理 4，存在自然数 m 使 \((\mathrm{A}_{\mathrm{M}})_{s}\) 同构于 \(M^{m}\) 的一个子模。由于 \(M^{m}\) 是半单 \(A_{M}\)-模，\((\mathrm{A}_{\mathrm{M}})_{s}\) 同构于 \(M^{m}\) 的一个直因子子模，故 M 是生成 \(A_{M}\)-模（定理 1）。
\item
  设 A 是主理想整环，P 是有限生成 A-模。存在整数 \(n \geqslant 1\) 与 A 的递增理想序列 \((\mathfrak{a}_{i})_{1 \leqslant i \leqslant n}\)，使 P 同构于诸 \(A/\mathfrak{a}_{i}\) 的直和（VII, §4, No.~4, 第 19 页，定理 2）；P 的零化理想 a 等于 \(a_{1}\)。于是 P 是环 A/a 上的生成模。若 P 不是挠模，则我们有 a = 0，而 P 是生成 A-模。
\end{enumerate}

引理 1. —— 设 A 是交换环，M 是有限生成 A-模，Ann(M) 是它的零化子。设 \(\mathfrak{a}\) 是 A 的一个理想。下列性质等价：

\begin{enumerate}
\def\labelenumi{(\roman{enumi})}
\item
  aM = M.
\item
  Ann(M) + a = A.
\item
  存在一个元素 \(a\)（属于 \(\mathfrak{a}\)），使得 \(am = m\) 对每个 \(m \in \mathrm{M}\) 成立。 (i) \(\Rightarrow\) (ii)：设 \((x_1, \ldots, x_n)\) 是 A-模 M 的一个生成族。若 \(\mathfrak{aM} = \mathrm{M}\)，则每个 \(x_i\) 都可写成 \(\sum_{j=1}^{n} c_{ij} x_j\)，其中 \(c_{ij}\) 属于 \(\mathfrak{a}\)。把矩阵 \((c_{ij})\) 记作 \(C\)，把以 \(x_1, \ldots, x_n\) 为元素的列矩阵记作 \(X\)。我们有 \((I_n - C)X = 0\)。设 \(d\) 为行列式，\(V\) 为矩阵 \(I_n - C\) 的余子式矩阵。由 III, §8, No.~6, 第 532 页的公式 (26)，我们有 \(dX = {}^t V(I_n - C)X = 0\)，从而 \(d \in \operatorname{Ann}(\mathrm{M})\)。另一方面，由于 \(c_{ij}\) 属于 \(\mathfrak{a}\)，我们有 \(d \equiv 1 (\bmod \mathfrak{a})\)（III, §8, No.~5, 第 529 页，(18)），因而 \(1 \in \operatorname{Ann}(\mathrm{M}) + \mathfrak{a}\)。
\item
  \(\Rightarrow\) (iii)：设 (ii) 成立，则存在 \(a\in \mathfrak{a}\) 与 \(b\in \mathrm{Ann}(\mathrm{M})\) 使 \(a + b = 1\)。于是 \(am = m\) 对每个 \(m\in \mathrm{M}\) 成立。
\item
  \(\Rightarrow\) (i)：这是明显的。
\end{enumerate}

命题 3. —— 设 A 是交换环。每个有限生成的忠实投射 A-模都是生成模。更一般地，有限生成投射 A-模 P 是生成 \(A_{P}\)-模。

设 P 是有限生成投射 A-模。我们有 \(\tau(\mathrm{P})\mathrm{P}=\mathrm{P}\)（VIII, 第 80 页）。

若 A-模 P 是忠实的，则理想 \(\tau(\mathrm{P})\) 等于 A（引理 1），且 A-模 P 是生成模（定理 1）；这就得到第一条断言。第二条断言由下面的引理 2 得出。

引理 2. —— 设 A 是一个环，M 是投射 A-模。则 \(\mathrm{A_M}\)-模 M 是投射的。

设 \((x_{i})_{i\in\mathrm{I}}\) 是 A-模 M 的一个生成族。存在 A-模 M 上的一族线性型 \((x_{i}^{*})_{i\in\mathrm{I}}\)，使对每个 \(x\in\mathrm{M}\)，族 \((\langle x,x_{i}^{*}\rangle)_{i\in\mathrm{I}}\) 有有限支集，且 \(x=\sum_{i\in\mathrm{I}}\langle x,x_{i}^{*}\rangle x_{i}\)（II, §2, No.~6, 第 238 页，命题 12）。对每个 \(i\in\mathrm{I}\)，映射 \(x\mapsto\langle x,x_{i}^{*}\rangle_{\mathrm{M}}\) 是 \(A_{M}\)-模 M 上的线性型，且 \(x=\sum_{i\in\mathrm{I}}\langle x,x_{i}^{*}\rangle_{\mathrm{M}}x_{i}\) 对每个 \(x\in\mathrm{M}\) 成立。由同前所引，M 是投射 \(A_{M}\)-模。

\subsection{生成模的双交换子}

定理 2. —— 生成模是平衡的。

设 A 是一个环，M 是生成 A-模；按定义，存在整数 \(n \geqslant 0\)、M 的元素 \(x_{1}, \ldots, x_{n}\) 及 \(x_{1}^{*}, \ldots, x_{n}^{*}\)（M 的对偶 \(M^{*}\) 的元素），满足 \(\sum_{i=1}^{n} \langle x_{i}, x_{i}^{*} \rangle = 1\)。回忆（II, §4, No.~2, 第 271 页），我们用下述公式定义群同态 \(\theta : M^{*} \otimes_{A} M \to \operatorname{End}_{A}(M)\)：\(\theta(x^{*} \otimes y)(x) = \langle x, x^{*} \rangle y\)。若 u 是 M 的双交换子的一个元素，则它与 \(\operatorname{End}_{A}(M)\) 交换；因此，对每个 \(y \in M\)，我们有

\[
\begin{array}{l} u (y) = u \Big (\sum_ {i = 1} ^ {n} \langle x _ {i}, x _ {i} ^ {*} \rangle y \Big) = \sum_ {i = 1} ^ {n} u (\theta (x _ {i} ^ {*} \otimes y) (x _ {i})) \\ = \sum_ {i = 1} ^ {n} \theta (x _ {i} ^ {*} \otimes y) (u (x _ {i})) = \Big (\sum_ {i = 1} ^ {n} \langle u (x _ {i}), x _ {i} ^ {*} \rangle \Big) y. \end{array}
\]

因此，u 属于 \(A_{M}\)，从而 M 是平衡的。

推论 1. —— 自由模是平衡的。

当模为零时这是明显的，而非零自由模是生成模（VIII, 第 81 页，例 2）。

推论 2. —— 设 A 是一个环，n 是一个整数 \(\geqslant 0\)。\(\mathbf{M}_{n}(\mathrm{A})\) 的中心由元素属于 A 的中心的纯量矩阵组成。我们把 \(A^{n}\) 视为左 \(\mathbf{M}_{n}(\mathrm{A})\)-模（II, §10, No.~7, 第 349 页）。这个模的自同态是映射 \(x \mapsto xa\)，其中 a 取遍 A。

设 M 是右 A-模 \(A_{d}^{n}\)。由推论 1，它是平衡的。于是 \(A_{M}\)、\(A_{M}^{\prime}\) 与 \(A_{M}^{\prime\prime}\) 的中心重合。再由 \(A_{M}\) 与 A 等同且 \(A_{M}^{\prime}\) 与 \(\mathbf{M}_{n}(\mathbf{A})\) 等同，即得推论 2。

注. —— 设 M 是一个 A-模。若 \(A_{M}\)-模 M 是生成模，则我们有 \(A_{M} = A_{M}^{\prime\prime}\)（这由把定理 2 应用于 \(A_{M}\)-模 M 得到），从而 M 是平衡 A-模。

推论 3. —— 交换环上的每个有限生成投射模都是平衡的。

事实上，由 VIII, 第 82 页的命题 3，有限生成投射模 M 是生成 \(A_{M}\)-模。因此本推论由上面的注得出。

推论 4. —— 主理想整环上的每个有限生成模都是平衡的。

事实上，主理想整环 A 上的有限生成模 M 是生成 \(A_{M}\)-模（VIII, 第 81 页，例 4）。

推论 5. —— 设 K 是一个交换域（域），V 是 K 上一个有限维向量空间，u 和 v 是 V 的自同态。下列性质等价：

\begin{enumerate}
\def\labelenumi{(\roman{enumi})}
\item
  存在一个多项式 \(\mathrm{P}\)，它属于 \(\mathrm{K}[\mathrm{X}]\)，使得 \(v = \mathrm{P}(u)\)。
\item
  自同态 v 与 V 的每个与 u 交换的自同态交换。
\end{enumerate}

取 A 为环 K{[}X{]}，M 为由 V 和 u 得到的 K{[}X{]}-模（VII, §5, No.~1, 第 28 页）。断言 (i) 表示 \(v \in K[X]_{M}\)，(ii) 表示 \(v \in K[X]_{M}'\)。因此推论 5 是推论 4 的一个特例。

命题 4. —— 反模为有限生成的半单模是平衡的。

这由 VIII, 第 81 页的例 3 及该注得出。

推论 1. —— 设 \((\mathrm{S}_{i})_{i\in\mathrm{I}}\) 是一个由两两不同构的单 A-模组成的有限族。对 \(i\in I\)，用 \(D_{i}\) 表示 \(S_{i}\) 的自同态体之反环。假设对每个 \(i\in I\)，\(D_{i}\)-向量空间 \(S_{i}\) 是有限维的。则映射 \(a \mapsto (a_{\mathrm{S}_i})_{i \in \mathrm{I}}\)（从 A 到 \(\prod_{i \in \mathrm{I}} \mathrm{End}_{\mathrm{D}_i}(\mathrm{S}_i)\)）是满射。

考虑 A-模 \(M = \prod_{i \in I} S_i\)。由于 I 有限，我们又有 \(M = \bigoplus_{i \in I} S_i\)，且 \(S_i\) 在 M 中的像是 M 的 \(S_i\) 型同型分量。因此 A-模 M 的自同态是映射 \((s_i) \mapsto (s_i d_i)\)，其中 \((d_i)_{i \in I}\) 取遍 \(\prod_{i \in I} D_i\)（VIII, 第 66 页，命题 5）。由于 I 有限且对每个 \(i \in I\)，右 \(D_i\)-向量空间 \(S_i\) 是有限维的，M 的反模是有限生成的。由命题 4，A-模 M 是平衡的。而 A-模 M 的双交换子由 \(\text{End}_{\mathbf{Z}}(\text{M})\) 中形如 \(\prod_{i \in I} u_i\) 的元素组成，其中 \((u_i) \in \prod_{i \in I} \text{End}_{\text{D}_i}(\text{S}_i)\)（VIII, 第 78 页，命题 1），因为 \(\text{End}_{\text{D}_i}(\text{S}_i)\) 是 \(S_i\) 的双交换子。本推论由此得出。

特别地，当 A 是交换域（域）K 上的代数且每个 \(S_{i}\) 是作为 K-向量空间有限维的单 A-模时，本推论适用：事实上，\(D_{i}\) 这时含诸位似 \(\alpha_{S_{i}}\)，其中 \(\alpha\) 取遍 K，且 \(S_{i}\) 在 \(D_{i}\) 上有限维，因为它在 K 上如此。

推论 2（伯恩赛德定理）. —— 设 A 是代数闭交换域（域）K 上的代数，S 是一个作为 K-向量空间有限维的单 A-模。则映射 \(a \mapsto a_{\mathrm{S}}\)（从 A 到 \(\operatorname{End}_{\mathrm{K}}(\mathrm{S})\)）是满射。

事实上，A-模 S 的自同态体由诸位似 \(\alpha_{S}\)（\(\alpha \in K\)）组成（VIII, 第 47 页，定理 1）。于是我们对单 A-模 S 应用推论 1。

\subsection{半单模的反模}

设 A 是一个环。把单 A-模的类组成的集合记作 S。对每个 \(\lambda \in S\)，选取一个单 A-模 \(S_{\lambda}\)，其类为 \(\lambda\)，并把 \(S_{\lambda}\) 的自同态体之反环记作 \(D_{\lambda}\)。我们把 \(S_{\lambda}\) 视为 \((A, D_{\lambda})\)-双模。

设 \(M\) 是一个半单 A-模，\(B\) 是 \(M\) 的自同态环。把 \(M\) 的双交换子记作 \(C\)。对每个 \(\lambda \in \mathcal{S}\)，把左 \((D_{\lambda}, B)\)-双模 \(\mathrm{Hom}_A(S_{\lambda}, M)\) 记作 \(V_{\lambda}\)。最后，把 A-模 \(M\) 的支集记作 \(\mathcal{S}_M\)（VIII, 第 66 页）；它也是元素 \(\lambda\)（属于 \(\mathcal{S}\)）组成的集合，这些元素使 \(V_{\lambda}\) 非零。

注 1. —— A-模 M 的典范描述 \(\alpha_{M}\) 是左 (A, B)-双模的同构。由 VIII, 第 71 页命题 9 的推论，映射 \(f \mapsto (\operatorname{Hom}(1_{\mathrm{S}_{\lambda}}, f))_{\lambda \in \mathcal{S}_{\mathrm{M}}}\)（从 B 到 \(\prod_{\lambda \in \mathcal{S}_{\mathrm{M}}} \operatorname{End}_{\mathrm{D}_{\lambda}}(\mathrm{V}_{\lambda})\)）是环同构。

命题 5. —— a) M 的反模是半单的。

\begin{enumerate}
\def\labelenumi{\alph{enumi})}
\setcounter{enumi}{1}
\item
  对每个 \(\lambda \in \mathcal{S}_{\mathrm{M}}\)，B-模 \(V_{\lambda}\) 是单模，其交换子等于 \((D_{\lambda})_{V_{\lambda}}\)。
\item
  映射 \(\lambda \mapsto \mathrm{cl}(\mathrm{V}_{\lambda})\) 是从 A-模 M 的支集到其反模的支集的双射。
\item
  对每个 \(\lambda \in \mathcal{S}_{\mathrm{M}}\)，B-子模 \(\mathrm{M}_{\lambda}\) 是 B-模 M 的 \(\mathrm{V}_{\lambda}\) 型同型分量，且 \(\mathrm{V}_{\lambda}\) 在 M 中的重数等于 \(\dim_{\mathrm{D}_{\lambda}}(\mathrm{S}_{\lambda})\)。
\item
  对 \(s \in S\)，把映射 \(\varphi \mapsto \varphi(s)\)（从 \(V_{\lambda} = \operatorname{Hom}_{\mathrm{A}}(\mathrm{S}_{\lambda}, \mathrm{M})\) 到 M）记作 \(\widetilde{s}\)。它是 B-线性的。所得到的映射 \(s \mapsto \widetilde{s}\)（从 \(S_{\lambda}\) 到 \(\operatorname{Hom}_{\mathrm{B}}(\mathrm{V}_{\lambda}, \mathrm{M})\)）是 \((\mathrm{A}, \mathrm{D}_{\lambda})\)-双模的同构。
\end{enumerate}

设 \(\lambda \in \mathcal{S}_{\mathrm{M}}\)。把环 \(\mathrm{End}_{\mathrm{D}_{\lambda}}(\mathrm{V}_{\lambda})\) 记作 \(\mathrm{E}_{\lambda}\)；由于 \(\mathrm{V}_{\lambda}\) 是非零 \(\mathrm{D}_{\lambda}\)-向量空间，它是单 \(\mathrm{E}_{\lambda}\)-模（VIII, 第 45 页，例 3），其交换子等于 \((\mathrm{D}_{\lambda})_{\mathrm{V}_{\lambda}}\)（VIII, 第 82 页，定理 2 的推论 1）。由于 \(\mathrm{E}_{\lambda}\) 是 B-模 \(\mathrm{V}_{\lambda}\) 的位似环（VIII, 第 71 页，命题 9 的推论），这就证明了 b)。

M 的典范描述 \(\alpha_{M}\) 定义了一个同构 \(\alpha_{\lambda}\)（从 \(V_{\lambda}\otimes_{D_{\lambda}^{\circ}}S_{\lambda}\) 到 \(M_{\lambda}\)）。由于 \(V_{\lambda}\) 是单 B-模，B-模 \(V_{\lambda}\otimes_{D_{\lambda}^{\circ}}S_{\lambda}\) 是 \(V_{\lambda}\) 型同型模（VIII, 第 61 页，命题 1）；因此 B-模 \(M_{\lambda}\) 亦然，这就证明了 a)。

由上面的注 1，存在 B 的元素 \(e_{\lambda}\)（\(\lambda\) 取遍 \(\mathcal{S}_{\mathrm{M}}\)），使得 \((e_{\lambda})_{\mathrm{V}_{\lambda}} = 1_{\mathrm{V}_{\lambda}}\)，且 \((e_{\lambda})_{\mathrm{V}_{\mu}} = 0\)（当 \(\mu \in \mathcal{S}_{\mathrm{M}}\) 而 \(\mu \neq \lambda\)）。因此单 B-模 \(V_{\lambda}\) 两两不同构，这就证明了 c) 及 d) 的第一条断言。B-模 \(M_{\lambda}\) 同构于 \(V_{\lambda} \otimes_{D_{\lambda}^{\circ}} S_{\lambda}\)，故 \(\dim_{D_{\lambda}}(S_{\lambda})\) 是 \(V_{\lambda}\) 在 M 中的重数（II, §3, No.~7, 第 255 页，推论 1）。

映射 \(\sum_{\lambda\in\mathcal{S}_{M}}\alpha_{\lambda}\)（从 \(\bigoplus_{\lambda\in\mathcal{S}_{M}}V_{\lambda}\otimes_{D_{\lambda}^{\circ}}S_{\lambda}\) 到 M）给出半单 B-模 M 的一个描述（VIII, 第 69 页，定义 5）。由 VIII, 第 70 页命题 8, b)，对每个 \(\lambda\in\mathcal{S}_{M}\)，e) 中所述的从 \(S_{\lambda}\) 到 \(\mathrm{Hom}_{\mathrm{B}}(\mathrm{V}_{\lambda},\mathrm{M})\) 的映射是双射且 \(D_{\lambda}\)-线性。由于它显然是 A-线性的，这就证明了 e)。

注 2. —— 由证明可知，映射

\[
\sum_ {\lambda \in \mathcal {S} _ {\mathrm{M}}} V _ {\lambda} \otimes_ {D _ {\lambda} ^ {\mathrm{o}}} S _ {\lambda} \longrightarrow M
\]

（由 M 的典范描述所诱导）是 M 的反模的一个描述。

命题 6. —— a) 视为 (A, B°)-双模时，M 是半单的。

\begin{enumerate}
\def\labelenumi{\alph{enumi})}
\setcounter{enumi}{1}
\item
  对每个 \(\lambda \in \mathcal{S}_{\mathrm{M}}\)，\(\mathrm{M}_{\lambda}\) 是 M 的单 (A, B°)-子双模。
\item
  对 M 的每个 \((\mathrm{A},\mathrm{B}^{\circ})\)-子双模 N，存在一个唯一的子集 \(\Lambda\)（\(\mathcal{S}_{\mathrm{M}}\) 的），使得 N 等于 \(\oplus_{\lambda \in \Lambda}\mathrm{M}_{\lambda}\)。
\end{enumerate}

设 \(\lambda\) 属于 \(\mathcal{S}_{\mathrm{M}}\)。左 A-模 \(S_{\lambda}\) 与右 B-模 \(V_{\lambda}\) 都是单模，且 \(D_{\lambda}\) 是 \(S_{\lambda}\) 的自同态体之反环。由 VIII, 第 63 页的推论 2，(A, B°)-双模 \(S_{\lambda} \otimes_{D_{\lambda}} V_{\lambda}\) 是单的，而与它同构的 \(M_{\lambda}\) 亦然。这就证明了 b)，a) 随之可得。

若 \(\lambda\) 与 \(\mu\) 在 \(S_{M}\) 中不同，则 \(M_{\lambda}\) 与 \(M_{\mu}\) 作为 A-模不同构，更不必说作为 \((A,B^{\circ})\)-双模了。断言 c) 由 VIII, 第 68 页的推论 2 得出。

命题 7. —— a) 对 M 的双交换子 C 的每个元素 c 及每个 \(\lambda\in\mathcal{S}_{M}\)，存在唯一的元素 \(c_{\lambda}\)（属于 \(\operatorname{End}_{\mathrm{D}_{\lambda}}(\mathrm{S}_{\lambda})\)），使得对每个 \(\varphi\in\operatorname{Hom}_{\mathrm{A}}(\mathrm{S}_{\lambda},\mathrm{M})\) 及每个 \(s\in S_{\lambda}\)，我们有 \(c\varphi(s)=\varphi(c_{\lambda}s)\)。

\begin{enumerate}
\def\labelenumi{\alph{enumi})}
\setcounter{enumi}{1}
\item
  给 \(S_{\lambda}\)（\(\lambda\) 取遍 \(S_{M}\)）赋以由 a) 定义的 C-模结构。则典范映射 \(\alpha_{M}\)（从 \(\bigoplus_{\lambda\in S_{M}}S_{\lambda}\otimes_{D_{\lambda}}V_{\lambda}\) 到 M）是 \((\mathrm{C},\mathrm{B}^{\circ})\)-双模的同构。
\item
  映射 \(c \mapsto (c_{\lambda})_{\lambda \in \mathcal{S}_{\mathrm{M}}}\) 是从 C 到 \(\prod_{\lambda \in \mathcal{S}_{\mathrm{M}}}\mathrm{End}_{\mathrm{D}_{\lambda}}(\mathrm{S}_{\lambda})\) 的同构。
\end{enumerate}

断言 a) 与 c) 由 VIII, 第 71 页命题 9 得出，因为典范映射 \(\alpha_{\mathrm{M}}\)（从 \(\bigoplus_{\lambda \in \mathcal{S}} (\mathrm{S}_{\lambda} \otimes_{\mathrm{D}_{\lambda}} \mathrm{W}_{\lambda})\) 到 M）给出 B-模 M 的一个描述（注 2）。此外，\(\alpha_{\mathrm{M}}\) 是 (C, B°)-线性的，这就证明了 b)。

注 3. —— 给 \(S_{\lambda}\)（对每个 \(\lambda \in S_{M}\)）赋以命题 7, a) 给出的 C-模结构。若在命题 5（VIII, 第 85 页）中把 A 换成 B、把 B 换成 C，则可看出：对每个 \(\lambda \in S_{M}\)，左 C-模 \(S_{\lambda}\) 是单模，其交换子为 \(D_{\lambda}\)；C-模 M 的 \(S_{\lambda}\) 型同型分量等于 \(M_{\lambda}\)；且映射 \(\lambda \mapsto \mathrm{cl}_{\mathrm{C}}(\mathrm{S}_{\lambda})\) 是从 A-模 M 的支集到 C-模 M 的支集的双射。最后注意，从 \(S_{\lambda}\) 到 M 的 A-线性映射与 C-线性映射相同，M 的 A-子模与 C-子模相同，且环 \(\operatorname{End}_{\mathrm{A}}(\mathrm{M})\) 与 \(\operatorname{End}_{\mathrm{C}}(\mathrm{M})\) 相等。

设 M 是一个半单 A-模。把 A-模 M 的双交换子 C 的中心记作 Z；它也是 M 的交换子 B 的中心。给 M 及 \(S_{\lambda}\)（\(\lambda \in S_{M}\)）赋以由 C-模结构经标量限制得到的 Z-模结构。对每个 \(\lambda \in S_{M}\)，把体 \(D_{\lambda}\) 的中心记作 \(Z_{\lambda}\)。

命题 8. —— a) 映射 \(z \mapsto (z_{\mathrm{S}_{\lambda}})_{\lambda \in \mathcal{S}_{\mathrm{M}}}\) 是从 Z 到诸体 \(\mathbf{Z}_{\lambda}\) 的乘积的同构。

\begin{enumerate}
\def\labelenumi{\alph{enumi})}
\setcounter{enumi}{1}
\item
  A-模 M 是同型的且非零，当且仅当 Z 是一个体。
\item
  设 \(\Lambda\) 是 \(\mathcal{S}_{\mathrm{M}}\) 的一个子集。用 \(e_{\Lambda}\) 表示 \(\mathrm{Z}\) 中唯一的元素，它满足：\((e_{\Lambda})_{\mathrm{S}_{\lambda}} = 1_{\mathrm{S}_{\lambda}}\)（对 \(\lambda \in \Lambda\)），且 \((e_{\Lambda})_{\mathrm{S}_{\lambda}} = 0\)（对 \(\lambda \in \mathcal{S}_{\mathrm{M}} - \Lambda\)）。我们有 \((e_{\Lambda})_{\mathrm{M}_{\lambda}} = 1_{\mathrm{M}_{\lambda}}\)（对 \(\lambda \in \Lambda\)）及 \((e_{\Lambda})_{\mathrm{M}_{\lambda}} = 0\)（对 \(\lambda \in \mathcal{S}_{\mathrm{M}} - \Lambda\)）。
\item
  若 M 的支集 \(S_{M}\) 有限，则映射 \(\Lambda \mapsto e_{\Lambda}Z\) 是从 \(S_{M}\) 的子集组成的集合到 Z 的理想组成的集合的双射，而映射 \(a \mapsto aM\) 是从 Z 的理想组成的集合到 M 的 \((A,B^{o})\)-子双模组成的集合的双射。这些双射都是有序集的同构。逆双射把 M 的 \((A,B^{o})\)-子双模 N 映到由 Z 中把 N 送入 M 的元素 z 组成的理想。
\end{enumerate}

对 \(\lambda \in \mathcal{S}_{\mathrm{M}}\)，\(Z_{\lambda}\) 是交换子 \(D_{\lambda}\) 与双交换子 \(C_{\lambda}\)（A-模 \(S_{\lambda}\) 的）的公共中心。由上面的命题 7, c)，映射 \(c \mapsto (c_{\lambda})_{\lambda \in \mathcal{S}_{\mathrm{M}}}\) 是从 C 到 \(\prod_{\lambda \in \mathcal{S}_{\mathrm{M}}} C_{\lambda}\) 的同构。限制到中心上，我们得到同构 \(z \mapsto (z_{S_{\lambda}})_{\lambda \in \mathcal{S}_{\mathrm{M}}}\)（从 Z 到 \(\prod_{\lambda \in \mathcal{S}_{\mathrm{M}}} Z_{\lambda}\)），由此得 a)。

环 \(\prod_{\lambda \in \mathcal{S}_{\mathrm{M}}}Z_{\lambda}\) 是一个体，当且仅当集合 \(\mathcal{S}_{\mathrm{M}}\) 只有一个元素，由此得 b)。

断言 c) 由命题 7, a) 得出。设 \(\mathcal{S}_{\mathrm{M}}\) 有限。由 a) 及 I, §8, No.~10, 第 109 页的命题 8，映射 \(\Lambda \mapsto e_{\Lambda}\mathrm{Z}\) 是从 \(\mathfrak{P}(\mathcal{S}_{\mathrm{M}})\) 到 Z 的理想组成的集合的有序集同构。设 \(\Lambda\) 是 \(\mathcal{S}_{\mathrm{M}}\) 的一个子集。由 c)，我们有关系式 \(e_{\Lambda}\mathrm{ZM} = e_{\Lambda}\mathrm{M} = \bigoplus_{\lambda \in \Lambda}\mathrm{M}_{\lambda}\)；由 VIII, 第 86 页的命题 6, c)，剩下只需描述逆双射。而 \(z \in \mathrm{Z}\) 把 M 送入 \(\bigoplus_{\lambda \in \Lambda}\mathrm{M}_{\lambda}\) 当且仅当 \(z = e_{\Lambda}z\)，即 \(z \in e_{\Lambda}\mathrm{Z}\)。

推论. —— 假设 A 是代数闭交换域（域）K 上的代数，且 M 是一个作为 K 上向量空间有限维的半单 A-模。对每个元素 \(\lambda\)（在 \(S_{M}\) 中），用 \(e_{\lambda}\) 表示 M 的投影算子，其像为 \(M_{\lambda}\)，核为 \(\oplus_{\lambda\neq\mu}M_{\mu}\)。则 \((e_{\lambda})_{\lambda\in\mathcal{S}_{M}}\) 是 K 上向量空间 Z 的一个基。

由于 M 是 K 上有限维向量空间，且是非零子模族 \((\mathrm{M}_{\lambda})_{\lambda\in\mathcal{S}_{\mathrm{M}}}\) 的直和，故集合 \(S_{M}\) 有限，且每个空间 \(S_{\lambda}\)（\(\lambda\in S_{M}\)）都在 K 上有限维。由于体 K 代数闭，我们有 \(D_{\lambda}=Z_{\lambda}=K\)（VIII, 第 47 页，定理 1），且映射 \(z\mapsto(z_{\mathrm{S}_{\lambda}})_{\lambda\in\mathcal{S}_{\mathrm{M}}}\) 是从 Z 到 \(K^{S_{M}}\) 的同构（命题 8, a)）。于是本推论由命题 8 的 c) 得出。

\subsection{稠密性定理}

定理 3（雅各布森）. —— 设 M 是一个半单 A-模，c 是加群 M 的一个自同态。则 c 属于 M 的双交换子 \(A_{M}^{\prime\prime}\) 当且仅当它满足下列条件：

\begin{enumerate}
\def\labelenumi{(\Alph{enumi})}
\setcounter{enumi}{3}
\tightlist
\item
  对每个有限子集 \(\mathrm{F}\)（\(\mathrm{M}\) 的），存在一个元素 \(a\)（属于 \(\mathrm{A}\)），使得 \(c\) 与 \(a_{\mathrm{M}}\) 在 \(\mathrm{F}\) 上相合。
\end{enumerate}

首先设 \(c\) 满足条件 (D)。设 \(u\) 是 \(\mathrm{A}_{\mathrm{M}}^{\prime}\) 的一个元素。设 \(x\) 是 \(\mathrm{M}\) 的一个元素，对子集 \(\mathrm{F} = \{x, u(x)\}\) 应用条件 (D)。存在一个元素 \(a\)（属于 \(\mathrm{A}\)），使得 \(c(x) = ax\) 且 \(c(u(x)) = au(x)\)，于是 \(u(c(x)) = u(ax) = au(x) = c(u(x))\)。由于 \(x\) 是任意的，我们有 \(cu = uc\)；又因这对一切 \(u\) 成立，我们有 \(c \in \mathrm{A}_{\mathrm{M}}^{\prime \prime}\)。

为证其逆，我们使用下面的引理。

引理 3. —— 设 M 是一个半单 A-模。设 B 是 A-模 M 的双交换子。则 M 的每个 A-子模都是 M 的 B-子模。

设 N 是 M 的一个 A-子模。由 VIII, 第 56 页的推论 2，存在 A-模 M 的一个以 N 为像的投影算子 p；它属于 \(A_{M}^{\prime}\)。由于对每个 \(b \in B\) 我们有关系式 pb = bp，我们得到 N 是 M 的 B-子模。

现在来完成定理 3 的证明。设 c 属于 \(A_{M}^{\prime\prime}\)，且设 \(F = \{x_{1}, \ldots, x_{n}\}\) 是 M 的一个有限子集。把元素 \((x_{1}, \ldots, x_{n})\)（\(M^{n}\) 的）记作 x。A-模 \(M^{n}\) 是半单的，且由 VIII, 第 79 页的命题 2，它的双交换子与 \(A_{M}^{\prime\prime}\)-模 \(M^{n}\) 的位似重合。由引理 3，\(M^{n}\) 的 A-子模 Ax 是一个 \(A_{M}^{\prime\prime}\)-子模（\(M^{n}\) 的）。设 \(a \in A\) 使得 \((cx_{1}, \ldots, cx_{n})\) 等于 ax。于是 c 与 \(a_{M}\) 在 \(\{x_{1}, \ldots, x_{n}\}\) 上相合，这就蕴含条件 (D)。

注. —— 把 M 的加群的自同态环记作 E。给 M 赋以离散拓扑；环 E 由从 M 到 M 的映射组成，我们可以给它赋以由 \(M^{M}\) 上乘积拓扑诱导的拓扑（“M 中单收敛拓扑”，一般拓扑，I, §2, No.~3, 第 31 页）。E 上的拓扑是豪斯多夫的，且与 E 的加群结构相容。对 E 中每个 f，映射 \(g \mapsto f \circ g\) 与 \(g \mapsto g \circ f\)（从 E 到 E）都是连续的。因此，E 的每个子集的交换子在 E 中闭。于是定理 3 蕴含 \(A_{M}^{\prime\prime}\) 是 \(A_{M}\) 在 E 中的闭包。

\subsection{在体论中的应用}

命题 9. —— 设 L 是一个体，E 是 \(\operatorname{End}_{\mathbf{Z}}(\mathrm{L})\) 的一个子环，它含映射 \(\gamma_{a}: x \mapsto ax\)（对每个 \(a \in L\)）。用 K 表示 L 中满足下列条件的元素 a 组成的集合：对 L 中每个 x 及 E 中每个 u，\(u(xa) = u(x)a\)；它是 L 的一个子体。

设 V 是右 K-向量空间 L 的一个有限维线性子空间，h 是一个从 V 到 L 的 K-线性映射。存在 E 的一个元素在 V 上与 h 相合。

我们把 L 视为左 E-模。由于 E 含诸左乘 \(\gamma_{a}\)，L 的每个 E-子模都是体 L 的左理想；因此 E-模 L 是单模。加群 L 的与诸 \(\gamma_{a}\) 交换的每个自同态都形如 \(\delta_{b}: x \mapsto xb\)，其中 b 属于 L。于是 \(b \mapsto \delta_{b}\) 是从 K 到 \(\operatorname{End}_{\operatorname{E}}(\operatorname{L})\) 的反环的一个同构，而后者是一个体。

因此 E-模 L 的双交换子 \(E''\) 由右 K-向量空间 L 的自同态组成。设 v 是 K-向量空间 L 的一个自同态，它在 V 上的限制与 h 相合；它是 \(E''\) 的元素。设 \((x_{i})_{i\in\mathrm{I}}\) 是 V 在 K 上的一个基；由定理 3（VIII, 第 88 页），存在 E 的元素 u，使得 \(u(x_{i}) = v(x_{i}) = h(x_{i})\)（对每个 \(i \in I\)）。由线性性推得，对 V 中每个 x 有 \(u(x) = h(x)\)。

推论. —— 设 L 是一个体。设 \(\Gamma\) 是体 L 的自同构群的一个子群，K 是 \(\Gamma\) 的不变量体。设 V 是 L 的一个右 K-线性子空间，它在 K 上有有限维 n。则存在元素 \(\sigma_{1},\ldots,\sigma_{n}\)（属于 \(\Gamma\)），具有下列性质：对每个从 V 到 L 的 K-线性映射 u，存在 L 的元素 \(a_{1},\ldots,a_{n}\)，使得对 V 中每个 x 有 \(u(x)=\sum_{i=1}^{n}a_{i}\sigma_{i}(x)\)。

用 E 表示形如 \(x \mapsto \Sigma_{\sigma \in \Gamma} a_{\sigma} \sigma(x)\) 的从 L 到 L 的映射组成的集合，其中 \((a_{\sigma})_{\sigma \in \Gamma}\) 是 L 的元素组成的有限支集族。对 L 中每个 a 我们有 \(\gamma_{a} \in E\)，且 E 是加群 L 的自同态环的一个子环。此外，体 K 由 L 中满足下列条件的元素 a 组成：对 L 中每个 x 及 E 中每个 u，\(u(xa) = u(x)a\)。

设 H 是左 L-向量空间 \(\operatorname{Hom}_{\mathrm{K}}(\mathrm{V},\mathrm{L})\)；它具有维数 n。由命题 9，它由 \(\Gamma\) 的元素在 V 上的限制生成。存在 n 个元素 \(\sigma_{1},\ldots,\sigma_{n}\)（属于 \(\Gamma\)），它们在 V 上的限制构成 H 在 L 上的一个基。本推论由此得出。

注. —— 当体 L 交换时，本推论化为阿廷定理（V, §10, No.~6, 第 65 页）。

\BourbakiExercises{习题}

\begin{enumerate}
\def\labelenumi{\arabic{enumi})}
\tightlist
\item
  设 M 是一个 A-模，N 是 M 的一个直因子子模。
\end{enumerate}

\begin{enumerate}
\def\labelenumi{\alph{enumi})}
\item
  证明 N 在 M 的双交换子 \(A_{M}^{\prime\prime}\) 下稳定。
\item
  证明每个从 N 到 M 的 A-线性映射都是 \(A_{M}^{\prime\prime}\)-线性的。
\item
  证明在 N 上的限制定义一个环同态 \(A_{M}^{\prime\prime} \to A_{N}^{\prime\prime}\)。
\item
  假设 M/N 是一些同构于 N 的商的子模之和。证明典范同态 \(A_{M}^{\prime\prime} \rightarrow A_{N}^{\prime\prime}\) 是单的。若 N 是平衡的，则 M 亦然。
\end{enumerate}

\begin{enumerate}
\def\labelenumi{\arabic{enumi})}
\setcounter{enumi}{1}
\tightlist
\item
  设 K 是一个交换域（域），A 是 \(\mathbf{M}_{3}(\mathrm{K})\) 的由形如下式的矩阵组成的子环
\end{enumerate}

\[
\left( \begin{array}{c c c} a & 0 & 0 \\ b & c & 0 \\ 0 & 0 & a \end{array} \right)
\]

其中 \(a, b, c\) 属于 K。设 M 是 A-模 \(\mathrm{K}^3\)；它是子模 \(\mathrm{N} = \mathrm{Ke}_1 + \mathrm{Ke}_2\) 与 \(\mathrm{P} = \mathrm{Ke}_3\) 的直和。

\begin{enumerate}
\def\labelenumi{\alph{enumi})}
\item
  证明 M 是平衡的，典范映射 \(A_{M}^{\prime\prime} \rightarrow A_{N}^{\prime\prime}\)（习题 1）是单的但不是满的，而典范映射 \(A_{M}^{\prime\prime} \rightarrow A_{P}^{\prime\prime}\) 是满的但不是单的。
\item
  由此给出环 B 上的模 Q 及其直因子子模 R 的一个例子，使典范映射 \(B_{Q}^{\prime\prime} \rightarrow B_{R}^{\prime\prime}\) 既非单又非满（取 \(B = A \times A\) 及 \(Q = M \times M\)）。
\end{enumerate}

\begin{enumerate}
\def\labelenumi{\arabic{enumi})}
\setcounter{enumi}{2}
\item
  设 B 是一个环，I 是 B 的一个双边理想；设 A 是 \(B \times B\) 的由偶 \((x, y)\)（满足 \(x - y \in I\)）组成的子环。设 \(B_{1}\)（相应地，\(B_{2}\)）是赋以由第一个（相应地，第二个）投影导出的 A-模结构的加群 B。证明 A-模 \(B_{1}\) 与 \(B_{2}\) 都是平衡的。在 I 满足什么条件时 A-模 \(B_{1} \oplus B_{2}\) 是平衡的？
\item
  设 A 与 B 是环，P 是一个 (A, B)-双模。
\end{enumerate}

\begin{enumerate}
\def\labelenumi{\alph{enumi})}
\item
  假设 B-模 P 是生成模且 A-模 P 是平衡的。证明对每个平衡 B-模 M，A-模 \(P \otimes_{B} M\) 是平衡的。
\item
  假设 A-模 P 是生成模，且映射 \(b \mapsto b_{P}\)（从 B 到 \(\operatorname{End}_{\mathrm{A}}(\mathrm{P})\)）是双射。证明对每个生成 B-模 M，A-模 \(P \otimes_{B} M\) 是生成模。
\end{enumerate}

\begin{enumerate}
\def\labelenumi{\arabic{enumi})}
\setcounter{enumi}{4}
\item
  设 A 是一个环，M 是一个 A-模，N 是 M 的一个子模。证明若 A-模 M/N 是生成模，则 A-模 M 亦然。用一个例子证明 N 可以是生成模而 M 不是。
\item
  设 \((\mathrm{A}_i)_{i\in \mathrm{I}}\) 是一族环，A 是它们的乘积。证明 A-模 \(\bigoplus \mathrm{A}_i\) 是投射的且忠实的。它是否总是生成模？
\item
  设 A 是一个环，M 是一个 A-模。证明若 M 是生成模，则右 A-模 \(\mathrm{M}^{*} = \mathrm{Hom}_{\mathrm{A}}(\mathrm{M}, \mathrm{A}_{s})\) 是生成模。用一个例子证明 \(M^{*}\) 可以是生成模而 M 不是（参见习题 6）。
\item
  设 A 是一个主理想整环，M 是一个 A-模。证明下列性质等价：
\end{enumerate}

\begin{enumerate}
\def\labelenumi{(\roman{enumi})}
\item
  A-模 M 是生成模。
\item
  M 的对偶非零。
\item
  模 M 含有一个同构于 \(A_{s}\) 的直因子子模。
\end{enumerate}

\begin{enumerate}
\def\labelenumi{\arabic{enumi})}
\setcounter{enumi}{8}
\tightlist
\item
  设 \(D\) 是一个体，\(V\) 是 \(D\) 上一个右向量空间。我们把 \(V\) 视为环 \(A = \operatorname{End}_D(V)\) 上的一个左模。
\end{enumerate}

\begin{enumerate}
\def\labelenumi{\alph{enumi})}
\item
  证明 A-模 V 是投射的、有限生成的、忠实的且平衡的。证明 (D,A)-双模 \(\mathrm{Hom}_{\mathrm{A}}(\mathrm{V},\mathrm{A}_{s})\) 与 \(\mathrm{V}^{*}=\mathrm{Hom}_{\mathrm{D}}(\mathrm{V},\mathrm{D}_{d})\) 典范同构。
\item
  证明 A-模 V 是生成模当且仅当 D-向量空间 V 的维数有限。
\item
  A-模 V 的迹理想是什么？
\end{enumerate}

\begin{enumerate}
\def\labelenumi{\arabic{enumi})}
\setcounter{enumi}{9}
\tightlist
\item
  设 M 是一个 A-模。证明下列性质等价：
\end{enumerate}

\begin{enumerate}
\def\labelenumi{(\roman{enumi})}
\item
  A-模 M 是有限生成的。
\item
  对每个生成 A-模 P，存在一个自然数 n 及一个从 \(P^{n}\) 到 M 的满同态 f。
\end{enumerate}

\begin{enumerate}
\def\labelenumi{\arabic{enumi})}
\setcounter{enumi}{10}
\tightlist
\item
  设 \(\mathrm{P}\) 是一个有限生成投射 A-模；把它的迹理想记作 \(\tau(\mathrm{P})\)。\\
\end{enumerate}

\begin{enumerate}
\def\labelenumi{\alph{enumi})}
\item
  证明我们有 \(\tau(\mathrm{P})^2 = \tau(\mathrm{P})\) 及 \(\tau(\mathrm{P})\mathrm{P} = \mathrm{P}\)。
\item
  证明典范映射 \(\tau(\mathrm{P})\otimes_{\mathrm{A}}\mathrm{P}\to\mathrm{P}\) 是同构。
\end{enumerate}

¶ 12) 设 A 是一个环，P 是一个有限生成 A-模。证明下列性质等价：

\begin{enumerate}
\def\labelenumi{(\roman{enumi})}
\item
  A-模 P 是投射的且是生成模。
\item
  A-模 \(\mathrm{P}^{(\mathbf{N})}\) 与 \(\mathrm{A}_s^{(\mathbf{N})}\) 同构。
\end{enumerate}

（为证蕴含关系 \((\mathrm{i})\Rightarrow(\mathrm{ii})\)，用归纳法构造整数的严格递增序列 \((n_{k})\) 与 \((m_{k})\) 以及 A-模的满同态 \(f_{k}:A_{s}^{n_{k}}\to P^{m_{k}}\) 与 \(g_{k}:P^{m_{k+1}}\to A_{s}^{n_{k}}\)，使得 \(g_{k-1}\circ f_{k}\) 与 \(f_{k}\circ g_{k}\) 是典范投影。）

\begin{enumerate}
\def\labelenumi{\arabic{enumi})}
\setcounter{enumi}{12}
\item
  设 k 是一个环，n 是一个整数 \(\geqslant 2\)；用 A 表示 \(\mathbf{M}_{n}(k)\) 的由上三角矩阵组成的子环，用 M 表示右 A-模 \(k^{n}\)。证明 M 是投射的且有限生成的，但 M 的位似环在其双交换子中不稠密（参见 VIII, 第 88 页，定理 3 及 VIII, 第 88 页，注）。特别地，M 不是平衡的。
\item
  Z-模 Q 的双交换子是什么？证明 Q 是其双交换子上的单模，但 Z-模 Q 的位似环在其双交换子中不稠密。
\end{enumerate}

¶ 15) 设 M 是主理想整环上的一个挠模。证明 M 的位似环在其双交换子中稠密。

（化归为 M 是 A 的不可约元素 \(\pi\) 的 \(\pi\)-准素模的情形。然后证明 M 的每个有限子集 S 都含于一个直因子子模中，该子模是一个有限生成模与有限个不可分解可除模的直和。借助 VII, §2, 第 54 页习题 3 及 VII, §2, 第 56 页习题 8，对 S 的基数用归纳法。）

\begin{enumerate}
\def\labelenumi{\arabic{enumi})}
\setcounter{enumi}{15}
\tightlist
\item
  设 D 是一个体，V 是 D 上一个维数 \(\geqslant 2\) 的向量空间，A 是 \(\operatorname{End}_{\mathrm{D}}(\mathrm{V})\) 的一个子环。
\end{enumerate}

\begin{enumerate}
\def\labelenumi{\alph{enumi})}
\item
  假设对每个由四个元素组成的组 \((x, x', y, y')\)（其中 x 与 \(x'\) 线性无关），存在 A 的元素 u，使得 \(u(x) = y\) 且 \(u(x') = y'\)。证明 A-模 V 的交换子是 D（与 V 的位似环等同），且这个模的位似环在其双交换子中稠密。
\item
  给出环 \(A \subset \operatorname{End}_{\mathrm{D}}(\mathrm{V})\) 的一个例子，使 V 是单 A-模但其交换子异于 D（考虑一个含 D 的体 E、E 上的一个向量空间 V 以及环 \(A = \operatorname{End}_{\mathrm{E}}(\mathrm{V})\)）。
\end{enumerate}

\begin{enumerate}
\def\labelenumi{\arabic{enumi})}
\setcounter{enumi}{16}
\tightlist
\item
  设 K 是一个代数闭交换域（域），V 是 K 上一个有限维向量空间，G 是 \(\operatorname{Aut}(V)\) 的一个子幺半群，使得 K{[}G{]}-模 V 是单模。证明若标量 \(\operatorname{Tr}(g)\)（\(g \in G\)）组成的集合有限，则 G 有限（由伯恩赛德定理推出 G 含一个基 \((g_{i})_{i \in I}\)（\(\operatorname{End}(V)\) 的），并考虑映射 \(g \mapsto (\operatorname{Tr}(gg_{i}))\)（从 G 到 \(K^{I}\)））。
\end{enumerate}

¶ 18) 若群 G 的每个具有有限生成族的子群都有限，则称 G 为局部有限的。

\begin{enumerate}
\def\labelenumi{\alph{enumi})}
\item
  设 G 是一个群，H 是 G 的一个正规子群。证明若 H 与 G/H 都局部有限，则 G 亦然。
\item
  由 a) 推出，一个其元素均为有限阶的可解群是局部有限的。
\item
  设 K 是一个交换域（域），V 是 K 上一个有限维向量空间，G 是 Aut(V) 的一个子群，它有一个由有限阶元素组成的有限生成族。证明 G 的元素的阶是有界的。
\end{enumerate}

（通过考虑由 G 的一个有限生成族的系数生成的子域 K，化归为 K 是素域 P 的纯超越扩张的情形。然后证明：对 \(\mathbf{GL}_{n}(\mathrm{K})\) 的每个有限阶元素 s，我们有 \(\varphi(s) \leqslant n\)（当 P = Q），且 \(s \leqslant p^{n} - 1\)（当 \(P = F_{p}\)）。）

\begin{enumerate}
\def\labelenumi{\alph{enumi})}
\setcounter{enumi}{3}
\tightlist
\item
  证明 Aut(V) 的每个其元素均为有限阶的子群 H 都是局部有限的。\(^{(1)}\)
\end{enumerate}

（化归为 K 代数闭的情形，并对 \(\dim(V)\) 用归纳法。设 G 是 H 的一个具有有限生成族的子群。若 V 是单 K{[}G{]}-模，应用习题 17；若 V 含一个在 G 下稳定的非平凡子空间 W，考虑 G 在 \(\operatorname{Aut}(V) \times \operatorname{Aut}(V/W)\) 中的像，并应用 a) 与 b)。）

¶ 19) 设 A 是一个环，使得 A-模 \(A_{s}\) 是一族有限的左理想 \((\mathfrak{I}_{i})_{1\leqslant i\leqslant n}\)（作为 A-模彼此同构）的直和。

\begin{enumerate}
\def\labelenumi{\alph{enumi})}
\item
  证明在 A 中存在一族元素 \(e_{i,j}\)（\(1 \leqslant i \leqslant n\)，\(1 \leqslant j \leqslant n\)），使得 \(e_{i,j}e_{h,k} = \delta_{j,h}e_{i,k}\)，其中 \(\delta_{j,h}\) 表示克罗内克符号，且 \(I_i = Ae_{i,i}\)（参见 VIII, 第 39 页，习题 5 及 I, §8, 第 173 页，习题 12）。反之，若在 A 中存在一族 \(n^2\) 个元素 \(e_{i,j}\)，使得 \(e_{i,j}e_{h,k} = \delta_{j,h}e_{i,k}\) 且 \(1 = \sum_{i=1}^{n} e_{i,i}\)，则诸左理想 \(Ae_{i,i}\) 作为 A-模同构（参见 VIII, 第 39 页，习题 5）。此外，若 B 是 A 中与所有 \(e_{i,j}\) 交换的元素组成的子环，则 A 同构于矩阵环 \(\mathbf{M}_n(\mathbf{B})\)，且 B 同构于诸 A-模 \(Ae_{i,i}\) 中每一个的交换子之反环（利用 VIII, 第 78 页及 VIII, 第 65 页，推论）。
\item
  设 M 是一个 A-模。证明诸 \(N_{i} = e_{i,i}M\) 是 B-模，且 M 作为 B-模是诸 \(N_{i}\) 的直和；此外，诸 B-模 \(N_{i}\) 两两同构（考虑映射 \(x \mapsto e_{i,1}x\) 与 \(y \mapsto e_{1,i}y\)），且诸 \(N_{i}\) 的零化子等于 B 与 M 在 A 中的零化子之交。反之，对每个 B-模 N，在 n 个同构于 N 的 B-模的直和 M 上定义一个 A-模结构，使 \(e_{1,1}M\) 同构于 N。
\item
  对 \(N_{1}\) 的每个 B-子模 P，指定 M 的 A-子模 \(\sum_{i=1}^{n} e_{i,1}P\)。证明这定义了从 \(N_{1}\) 的 B-子模组成的集合到 M 的 A-子模组成的集合的双射，且它对包含关系严格递增。
\end{enumerate}

\section{模与代数的森田等价}

在本节中，\(k\) 表示一个交换环。

\subsection{交换子与对偶}

设 A 与 B 是 k-代数。回忆（III, §4, No.~3, 第 466 页）代数 A 与 B 上的双模是这样的 (A,B)-双模 P：由 A-模结构与 B-模结构导出的两个 k-模结构在它上面相合。为避免歧义，我们说 P 是 \((\mathrm{A},\mathrm{B})_{k}\)-双模。设 P 是一个 \((\mathrm{A},\mathrm{B})_{k}\)-双模。我们用 \(P^{*}\) 表示 P 的底层左 A-模的对偶 \(\operatorname{Hom}_{\mathrm{A}}(\mathrm{P},\mathrm{A})\)。它是一个 \((\mathrm{B},\mathrm{A})_{k}\)-双模（II, §1, No.~14, 第 226 页）；对 \(a \in A\)、\(b \in B\)、\(x \in P\) 及 \(x^{*} \in P^{*}\)，我们有

\[
\langle x, b x ^ {*} a \rangle = \langle x b, x ^ {*} \rangle a.\tag{1}
\]

我们用 \(_{s}A_{d}\) 表示视为 \((\mathrm{A},\mathrm{A})_{k}\)-双模（同前所引）的代数 A，并用 \(\Lambda:P\otimes_{B}P^{*}\to_{s}A_{d}\) 表示由下式确定的 \((\mathrm{A},\mathrm{A})_{k}\)-双模同态

\[
\Lambda (x \otimes x ^ {*}) = \langle x, x ^ {*} \rangle\tag{2}
\]

其中 \(x \in \mathrm{P}\) 且 \(x^{*} \in \mathrm{P}^{*}\) 。我们用 \(\widetilde{\mathrm{P}}\) 表示对偶 \(\mathrm{Hom}_{\mathrm{B}}(\mathrm{P},\mathrm{B})\)（\(\mathrm{P}\) 的底层右 B-模的）；它是一个 \((\mathrm{B},\mathrm{A})_{k}\)-双模。我们用 \(\widetilde{\Lambda}: \widetilde{\mathrm{P}} \otimes_{\mathrm{A}} \mathrm{P} \to {}_{s}\mathrm{B}_{d}\) 表示由下式确定的 \((\mathrm{B},\mathrm{B})_{k}\)-双模同态

\[
\widetilde {\Lambda} (\widetilde {x} \otimes x) = \langle \widetilde {x}, x \rangle\tag{3}
\]

其中 \(x \in P\)，\(\widetilde{x} \in \widetilde{P}\)。

现在，设映射 \(b \mapsto b_{\mathrm{P}}\) 是从 B 到 \(\operatorname{End}_{\mathrm{A}}(\mathrm{P})\) 的双射；这时它是从 B 到 \(\operatorname{End}_{\mathrm{A}}(\mathrm{P})\) 的反代数的同构。\(\mathbf{Z}\)-模的一个典范同态（从 \(\mathrm{P}^* \otimes_{\mathrm{A}} \mathrm{P}\) 到 \(\operatorname{End}_{\mathrm{A}}(\mathrm{P})\)；II, §4, No.~2, 第 271 页）于是确定了一个 \(\mathbf{Z}\)-模同态 \(\Theta: \mathrm{P}^* \otimes_{\mathrm{A}} \mathrm{P} \to \mathrm{B}\)，它由下式定义

\[
x \Theta (x ^ {*} \otimes y) = \langle x, x ^ {*} \rangle y\tag{4}
\]

其中 \(x, y \in \mathrm{P}\)，\(x^{*} \in \mathrm{P}^{*}\)。由 (1) 式，这个同态是 (B,B)-线性的，且我们有

\[
\Theta (x ^ {*} \otimes y) y ^ {*} = x ^ {*} \langle y, y ^ {*} \rangle\tag{5}
\]

其中 \(y \in \mathrm{P}\)，\(x^{*}, y^{*} \in \mathrm{P}^{*}\)。由 (4) 与 (5)，我们在 \((\mathrm{B}, \mathrm{B})_{k}\)-双模 \(\mathrm{P}^{*} \otimes_{\mathrm{A}} \mathrm{P}\) 中推得下列等式：

\[
(y ^ {*} \otimes x) \Theta (x ^ {*} \otimes y) = y ^ {*} \otimes \langle x, x ^ {*} \rangle y = y ^ {*} \langle x, x ^ {*} \rangle \otimes y = \Theta (y ^ {*} \otimes x) (x ^ {*} \otimes y)
\]

其中 \(x, y \in \mathrm{P}\)，\(x^{*}, y^{*} \in \mathrm{P}^{*}\)，从而

\[
s \Theta (t) = \Theta (s) t\tag{6}
\]

其中 \(s, t \in \mathrm{P}^* \otimes_{\mathrm{A}} \mathrm{P}\)。同样，对 \(x, y\)（在 \(\mathrm{P}\) 中）及 \(x^*, y^*\)（在 \(\mathrm{P}^*\) 中），我们由 (4) 与 (5) 在 \((\mathrm{A}, \mathrm{A})_k\)-双模 \(\mathrm{P} \otimes_{\mathrm{B}} \mathrm{P}^*\) 中推得下列等式：

\[
(x \otimes x ^ {*}) \langle y, y ^ {*} \rangle = x \otimes \Theta (x ^ {*} \otimes y) y ^ {*} = x \Theta (x ^ {*} \otimes y) \otimes y ^ {*} = \langle x, x ^ {*} \rangle (y \otimes y ^ {*}),
\]

从而

\[
u \Lambda (v) = \Lambda (u) v\tag{7}
\]

其中 \(u, v \in \mathrm{P} \otimes_{\mathrm{B}} \mathrm{P}^{*}\)。

对每个元素 \(x^{*}\)（属于 \(P^{*}\)），把从 P 到 B 的 B-线性映射 \(x \mapsto \Theta(x^{*} \otimes x)\) 记作 \(\sigma(x^{*})\)。这样我们就定义了一个映射 \(\sigma\)（从 \(P^{*}\) 到 \(\widetilde{P}\)），它是 (B, A)-线性的，且按定义满足

\[
\Theta (x ^ {*} \otimes y) = \langle \sigma (x ^ {*}), y \rangle\tag{8}
\]

其中 \(x^{*} \in P^{*}\)，\(y \in P\)。由 \(\widetilde{\Lambda}\) 的定义，我们因此有

\[
\Theta = \widetilde {\Lambda} \circ (\sigma \otimes 1 _ {\mathrm{P}}).\tag{9}
\]

设映射 \(a \mapsto a_{\mathrm{P}}\)（从 A 到 \(\operatorname{End}_{\mathrm{B}}(\mathrm{P})\)）是双射；这时它是一个代数同构。类似地，我们用下式定义一个 \((\mathrm{A}, \mathrm{A})_k\)-双模的同态 \(\widetilde{\Theta}: \mathrm{P} \otimes_{\mathrm{B}} \widetilde{\mathrm{P}} \to {}_s \mathrm{A}_d\)：

\[
\widetilde {\Theta} (x \otimes \widetilde {y}) y = x \langle \widetilde {y}, y \rangle\tag{10}
\]

其中 \(x, y \in \mathrm{P}\)，\(\widetilde{y} \in \widetilde{\mathrm{P}}\)。我们还用下式定义一个 \((\mathrm{B}, \mathrm{A})_k\)-双模的同态 \(\widetilde{\sigma}: \widetilde{\mathrm{P}} \to \mathrm{P}^*\)：

\[
\widetilde {\Theta} (x \otimes \widetilde {y}) = \langle x, \widetilde {\sigma} (\widetilde {y}) \rangle\tag{11}
\]

其中 \(x\in \mathrm{P},\widetilde{y}\in \widetilde{\mathrm{P}}\)。我们有

\[
\widetilde {\Theta} = \Lambda \circ (1 _ {P} \otimes \widetilde {\sigma}).\tag{12}
\]

命题 1. —— 设映射 \(b \mapsto b_{P}\)（从 B 到 \(\operatorname{End}_{\mathrm{A}}(\mathrm{P})\)）与映射 \(a \mapsto a_{P}\)（从 A 到 \(\operatorname{End}_{\mathrm{B}}(\mathrm{P})\)）都是双射。则 \(\sigma\) 与 \(\widetilde{\sigma}\) 是互逆的同构，且我们有关系式

\[
\Lambda = \widetilde {\Theta} \circ (1 _ {\mathrm{P}} \otimes \sigma),\tag{13}
\]

\[
\widetilde {\Lambda} = \Theta \circ (\widetilde {\sigma} \otimes 1 _ {\mathrm{P}}).\tag{14}
\]

对 \(x \in \mathrm{P}\)、\(x^{*} \in \mathrm{P}^{*}\) 及 \(y \in \mathrm{P}\)，由关系式 (4)、(8)、(10) 与 (11)，我们有

\[
\langle x, x ^ {*} \rangle y = x \Theta (x ^ {*} \otimes y) = x \langle \sigma (x ^ {*}), y \rangle = \widetilde {\Theta} (x \otimes \sigma (x ^ {*})) y = \langle x, \widetilde {\sigma} (\sigma (x ^ {*})) \rangle y.\tag{15}
\]

同样，对 \(x \in P\)、\(\widetilde{y} \in \widetilde{P}\) 及 \(y \in P\)，我们有

\[
x \langle \widetilde {y}, y \rangle = \widetilde {\Theta} (x \otimes \widetilde {y}) y = \langle x, \widetilde {\sigma} (\widetilde {y}) \rangle y = x \Theta (\widetilde {\sigma} (\widetilde {y}) \otimes y) = x \langle \sigma (\widetilde {\sigma} (\widetilde {y})), y \rangle .\tag{16}
\]

由这些假设，P 作为 A-模与作为 B-模都是忠实的。由关系式 (15) 与 (16)，我们分别推得 \(\widetilde{\sigma} \circ \sigma = 1_{P^{*}}\) 与 \(\sigma \circ \widetilde{\sigma} = 1_{\widetilde{P}}\)。关系式 (13) 与 (14) 于是分别由 (12) 与 (9) 得出。

注. —— 1) 设映射 \(b \mapsto b_{\mathrm{P}}\)（从 B 到 \(\operatorname{End}_{\mathrm{A}}(\mathrm{P})\)）是双射。则

\begin{enumerate}
\def\labelenumi{\alph{enumi})}
\item
  B-模 P 可与 P 的反模等同；因此它是忠实的且平衡的；
\item
  映射 \(a \mapsto a_{P}\)（从 A 到 \(\operatorname{End}_{\mathrm{B}}(\mathrm{P})\)）是双射，当且仅当 A-模 P 是忠实的且平衡的。
\end{enumerate}

\begin{enumerate}
\def\labelenumi{\arabic{enumi})}
\setcounter{enumi}{1}
\tightlist
\item
  在命题 1 的假设下，A-模 P 是平衡的；由于环 \(A_{P}\) 与 \(A_{P}^{\prime}\) 有相同的中心（VIII, 第 77 页），存在一个同构 \(\varphi\)（从环 A 的中心 \(Z(A)\) 到环 B 的中心 \(Z(B)\)），它由关系式 \(\varphi(z)_{\mathrm{P}} = z_{\mathrm{P}}\)（对 \(z \in \mathrm{Z}(A)\)）确定。此外，\((\mathrm{A}, \mathrm{B})_{k}\)-双模 P 的自同态是诸位似 \(z_{P}\)，其中 z 取遍 \(Z(A)\)；\((\mathrm{A}, \mathrm{B})_{k}\)-双模 P 的自同构是诸位似 \(z_{P}\)，其中 z 在 \(Z(A)\) 中可逆。
\end{enumerate}
